In this chapter, is presented the general proximal optimal consensus problem with the respected algorithm to be solved. Furthermore, the NOCGPROX variables are introduced to regulate the problem. In addition, a comprehensive proof for the validity of the problem is proposed.

\section{Optimal Proximal Consensus Problem}

    For a multi-agent system which is consisted $N$ agents with output $x_i \in \mathbf{R}^M$ and cost functions $f_i: \mathbf{R}^M \to \mathbf{R}$, $i \in \mathcal{V}$, the optimal proximal consensus problem which is introduced in \cite{9785920}, is achieved through a robust distributed control law for every agent, which ensures sublinear stability of the algorithm and 
    \begin{equation*}
        \lim_{t \to \infty} ||x(t) - x^*|| = 0
    \end{equation*}
    \begin{equation*}
        \lim_{t \to \infty} ||x^{r+1} - x^r|| = 0
    \end{equation*}
    \begin{equation*}
        \lim_{t \to \infty} ||\mu^{r+1} - \mu^r|| = 0
    \end{equation*}
    \begin{equation*}
        \lim_{t \to \infty} ||Ax^{r+1}|| = 0
    \end{equation*}
    where $x^* \in \mathbf{R}^Q$ is the optimal solution of (\ref{eq:2.5.6}). First and foremost, the main assumptions must be presented.
    \newpage
    \begin{assumption}\label{asm:1}
        The function $f(x)$ is differentiable and has Lipschitz continuous gradient, i.e.
        \begin{equation*}
            || \nabla f(x) - \nabla f(y) || \leq \mathrm{L} ||x - y||, \hspace{2.0mm} \forall x,y \in \mathbf{R}^Q
        \end{equation*}
        Further, it is assumed that $A^TA + B^TB \succeq I_Q$.
    \end{assumption}
    
    \begin{assumption}\label{asm:2}
        There exists a constant $\delta > 0$ such that
        \begin{equation*}
            \exists \hspace{2.0mm} \underline{f} > -\infty, \hspace{2.0mm} s.t. \hspace{2.0mm} f(x) + \frac{\delta}{2}||Ax||^2 \geq \underline{f}, \hspace{2.0mm} \forall x \in \mathbf{R}^Q
        \end{equation*}
        Without loss of generality and for the simplicity of notations, is going to be assumed that $\underline{f} = 0$.
    \end{assumption}
    
    \begin{assumption}\label{asm:3}
        The constraint $Ax = 0$ is feasible over $x \in \mathbf{R}^Q$.
    \end{assumption}
    
    \begin{assumption}\label{asm:4}
        The network topology is described by an undirected connected graph $\mathcal{G} = (\mathcal{V}, \mathcal{E})$, which is composed of $N$ agents (nodes) and is connected by $E$ edges. The agents communicate with each other only every $D$ sampled periods of time, i.e. every discrete points of time $t = DT$, where $T$ is the sampled period and $D$ a positive constant. In that time period, the agent $i$ will get samples of the output $x_j(DT)$ from his neighbors $j \in \mathcal{N}(i)$. Additionally, for a larger numerical value of $D$ the exchange of information between the agents will be conducted in lower frequency bands. In order to prevent communication delays during the exchange of the information, it is assumed that every measurement $x_j(DT)$ from the $j$ agents is obtainable by $i$ agent in time, $t_e = DT + d_{max}$, with $d_{max} \geq 0$ be the maximum permitted communication delay. Furthermore, it is assumed that the sampling period is selected sufficiently large, namely $D > d_{max}$.
    \end{assumption}

    With these basic assumptions for the algorithm, $B$ is chosen such that $B := |A|$, or \textit{signless incidence matrix}. Thus, it is known that $L_- = A^TA$, then $L_+ = B^TB$. Then it is certain that $A^TA + B^TB = 2\mathrm{D}$. So the minimum singular value of the matrix $A^TA$ is positive and is depended by the degrees of the graph. From equation (\ref{eq:2.5.8a}), the optimality condition is derived as
    \begin{equation}\label{eq:3.1.1}
        \nabla f(x^r) + A^T\mu^r + \beta A^TAx^{r+1} + \beta B^TB (x^{r+1} - x^r) = 0
    \end{equation}
    Rearranging the equation (\ref{eq:3.1.1}), the proximal update of the primal variable is
    \begin{equation*}
        \nabla f(x^r) + A^T\mu^r + \beta A^TAx^{r+1} + \beta B^TBx^{r+1} - \beta B^TBx^r = 0
    \end{equation*}
    \begin{equation*}
        \Rightarrow  \nabla f(x^r) + A^T\mu^r + \beta \left( A^TA + B^TB\right)x^{r+1} - \beta B^TBx^r = 0   
    \end{equation*}
    \begin{equation*}
        \Rightarrow    \nabla f(x^r) + A^T\mu^r + 2\beta \mathrm{D}x^{r+1} - \beta L_+x^r = 0
    \end{equation*}    
    \begin{equation*}
        \Rightarrow 2\beta \mathrm{D}x^{r+1} = \beta L_+x^r -A^T\mu^r -\nabla f(x^r)
    \end{equation*}
    Therefore, in vector form the update of the primal variable is
    \begin{equation}\label{eq:3.1.2}
        x^{r+1} = \frac{1}{2\beta}\mathrm{D}^{-1} \left[ \beta L_+x^r -A^T\mu^r -\nabla f(x^r)\right]
    \end{equation}
    or every agent is computing the position with
    \begin{equation} \label{eq:3.1.3}
        x_i^{r+1} = \frac{1}{2\beta d_i} \left[\beta \sum_{j \in \mathcal{N}(i)} x_j^r - \nabla f_i(x_i^r) - \beta\sum_{j \in \mathcal{N}(i)} \sum_{\ell = 1}^{r - 1} \left( x_i^{\ell} - x_j^{\ell}\right)\right]
    \end{equation}
    Thus, the primal update is computed by equation (\ref{eq:3.1.3}). The next step is the transformation of the problem with the proposed variables. 

\subsection{Definition of NOCGPROX Variables}

    In order to solve the proximal optimal consensus problem. there must be defined special variables that transform the described above optimization problem into a regulation one. Thus, it will be controllable with specific designed robust control laws. With regard to these variables, every agent will only exchange sampled output measurements in order to calculate them. The proposed regulation variables are called NOn-Convex Gradient PROXimal (NOCGPROX) variables of continuous time and defined as
    
    \begin{equation}\label{eq:3.1.4}
        q_i(t) := x_i(t) - \rho \left(t - \left\lfloor \frac{t}{T}\right\rfloor T \right) s_i \left( \left\lfloor \frac{t}{T}\right\rfloor T \right) - \left[ 1 - \rho \left( t - \left\lfloor \frac{t}{T}\right\rfloor T\right)\right] s_i \left(\left\lfloor \frac{t}{T}\right\rfloor T - T\right)
    \end{equation}
    where the discrete time variant vector $s_i$ is defined as:
    
    \begin{equation}\label{eq:3.1.5}
        \begin{align}
            s_i\left(\left \lfloor \frac{t}{T} \right \rfloor T \right) := x_i \left (\left \lfloor \frac{t}{T} \right \rfloor T \right) - \frac{1}{2\beta d_i}& \left[ \beta \sum_{j \in \mathcal{N}(i)}  x_j \left ( \left \lfloor \frac{t}{T} \right \rfloor T \right) - \left \nabla f_i \left( x_i \left( \left \lfloor \frac{t}{T} \right \rfloor T \right)\right) \right \\
            & - \left\beta \sum_{j \in \mathcal{N}(i)} \sum_{\ell = 1}^{\lfloor t/DT \rfloor} \left( x_i \left (\ell  T \right) - x_j \left ( \ell T \right) \right) \right]
        \end{align}
    \end{equation}
    for every $i \in \mathcal{V}$ with $\beta > 0$. The last term of equation (\ref{eq:3.1.5}) describes the exchange of information between the agents. The described $s_i$ derives from the selected algorithm and equation (\ref{eq:3.1.3}).
    
    The function $\rho:[0,\infty) \to [0,1]$ is a continuous function that is employed to generate the samples of the neighbors with the use of $s_i(kT)$ inside $q_i(t)$ in a continuous structure. A suitable function is presented in \cite{10125000}:
    \begin{equation}\label{eq:3.1.6}
        \rho(t) := \left \{ \begin{matrix}
            0, \hspace{50.0mm} t \in [0, d_{max})\\
            \frac{t - d_{max}}{T - d_{max}} - \frac{1}{2\pi}\sin \left( 2\pi \frac{t - d_{max}}{T - d_{max}}\right), \hspace{2.0mm} t \in [d_{max}, T] \\
            1, \hspace{50.0mm} t \in (T, \infty)
            \end{matrix} \right
    \end{equation}
    From definition, the function $\rho$ is increasing in its domain and is continuously differentiable for $t \geq 0$. Both derivates $\frac{d\rho}{dt}$ and $\frac{d^2\rho}{dt^2}$ are also differentiable in their domain. In time interval $[0, d_{max}]$ the auxiliary function is 0, thus NOCGPROX variables will use the neighbors' samples after $d_{max}$ time, when the last sample was received.

    \begin{remark}\label{rem:3.1}
        The function $\rho (t - \lfloor t/T \rfloor T)$ is not continuous for $t = kT$, $k \in \mathbf{N}$. On the contrary, the variables $q_i(t)$ are two times continuous differentiable as is $x_i(t)$. Therefore, for $t = kT$ it can be proven that the variables are continuous.
        \begin{equation*}
            \begin{align}
                &\lim_{t \to kT^-} q_i(t) = x_i(kT) - [1 - \rho (T)]s_i((k - 2)T) - \rho(T)s_i((k-1)T) = x_i(kT) - s_i((k-1)T)\\
                &\lim_{t \to kT^+} q_i(t) = x_i(kT) - [1 - \rho (0)]s_i((k-1)T) - \rho(0)s_i(kT) = x_i(kT) - s_i((k-1)T)
            \end{align}
        \end{equation*}
        Furthermore, hence $d^r \rho(t)/dt^r|_{t = kT} = 0$, for $r = 1,2,...$, the variables $q_i$ are continuous with continuous derivates with $q_i^{(r)}(kT) = x_i^{(r)}(kT)$
    \end{remark}

    \begin{remark}\label{rem:3.2}
        For $t \in (kT, (k+1)T)$, the variables have value $q_i(t) = x_i(t) - [1-\rho(t - kT)]s_i((k-1)T) - \rho(t - kT)s_i(kT)$. From definition of $\rho(t)$ in (\ref{eq:3.1.6}), $\rho(t - kT) = 0, \forall t \in (kT, kT + d_{max})$. As a result, the measurement $s_i(kT)$ needs only $t \geq kT + d_{max}$ for the computation of $q_i(t)$. So, the control law which will be applied in $q_i(t)$, will require the output measurements of the neighbors $x_j(kT), j \in \mathcal{N}(i)$, which is going to be inserted into the equation through $s_i(kT)$, after time $kT + d_{max}$. As a result, the communication between the agents will be protected from disturbances and delays.
    \end{remark}

    From the definition of the variables and the remark \ref{rem:3.1}, it is clearly that the variables $q_i(t)$ are actually position error of the agents. Thus, the new algorithm will be changed.
    \begin{equation*}
        q_i^{r+1} = x_i^{r+1} - s_i^r \Rightarrow
    \end{equation*}
    \begin{equation}\label{eq:3.1.7}
        x_i^{r+1} = q_i^{r+1} + s_i^r
    \end{equation}

\newpage
\subsection{Gradient Based Regulation Algorithm}
    
    The basic concept of integrating a discrete signal $s_i[k]$ in a continuous signal $q_i(t)$ for every agent was researched in \cite{9162430}. The next step is to change the algorithm of \cite{hong2016decomposing} in a discrete time form. From equation (\ref{eq:3.1.5}) the discrete time system is

    \begin{equation}\label{eq:3.1.8}
        \begin{align}
            q_i^{\nu + 1} = x_i^{\nu + 1} - x_i^{\nu} + \frac{1}{2\beta d_i}& \left[ \beta \sum_{j \in \mathcal{N}(i)} x_j^{\nu}  - \nabla f_i \left( x_i^{\nu}\right) \right\\
            & - \left\beta \sum_{j \in \mathcal{N}(i)} \sum_{\ell = 1}^{\lfloor \nu/D \rfloor} \left( x_i^{\ell} - x_j^{\ell}\right) \right], i = \{1,...,N\}
        \end{align}
    \end{equation}
    where $q:= col(q_1, q_2, ..., q_N) \in \mathbf{R}^Q$, $\nabla f(x) := col(\nabla f_1(x_1), \nabla f_2(x_2), ..., f_N(x_N))\in \mathbf{R}^Q$ and $x := col(x_1,  x_2, ..., x_N) \in \mathbf{R}^Q$. The equation (\ref{eq:2.5.8b}) updates the dual variable of the algorithm. In a vector discrete form it can be set transformed as:

    \begin{equation}\label{eq:3.1.9}
        \mu^{\nu D} := \beta \sum_{\ell = 1}^{\lfloor \nu/D \rfloor} Ax^{\ell}
    \end{equation}

    From the definition of (\ref{eq:3.1.9}) it is clearly that
    \begin{equation}\label{eq:3.1.10}
        \mu^{(\nu D + \ell)} = \left \{ \begin{matrix}
            \mu^{\nu D}, \hspace{20.0mm} \ell = 1,...,D-1 \\\\
            \mu^{\nu D} + \beta \sum_{\ell = 1}^{\lfloor \nu/D \rfloor} Ax^{\ell}, \hspace{5.0mm} \ell = D 
        \end{matrix} \right
    \end{equation}
    where $\nu \in \mathbf{N}$ and $\mu_D[\nu] = \mu^{\nu D}$. In addition, the dual variable will be updated only after $D$ sampled periods. Therefore, the discrete vector form of equation (\ref{eq:3.1.8}) will be
    \begin{equation}\label{eq:3.1.11}
        q^{\nu + 1} = x^{\nu + 1} - x^{\nu} + \frac{1}{2\beta} \mathrm{D}^{-1}\left[\beta L_+x^{\nu} - \nabla f(x^{\nu}) - \beta A^T\mu^{\nu D}\right]
    \end{equation}
    
    With the discrete form of all the signals of the system, in algorithm \ref{alg:reg-prox-gpda} is presented the regulated version of the Prox-GPDA algorithm. Every $D$ periods of time, the dual variable is updated. In addition, the output measurements of the agents are exchanging after $DT + d_{max}$ time, thus the new position is computed along with the control variables of the system. 
    
    \newpage
    \begin{algorithm}[H]
        \caption{The Regulated Proximal Primal Dual Algorithm (R-Prox-GPDA)}\label{alg:reg-prox-gpda}
        \textbf{Input}: $\mu^{\nu D} \in \mathbf{R}^{Q}$, $x^0 \in \mathbf{R}^Q$, $q^{\nu + 1} \in \mathbf{R}^Q$, $D, T, d_{max}, \beta$\\
        \textbf{Output}: $x^{\nu +1} \in \mathbf{R}^Q$, $s_i^{\nu} \in \mathbf{R}^Q$
        \begin{algorithmic}[1]
            \State $\nu \gets 1$;
            \State $\ell \gets 1$;\\
            \For {$i = 1:N$}{
                   
                \For{$j \in \mathcal{N}(i)$}{
                    $s_i^{\nu}$ =  $x_i^{\nu} -\frac{1}{2\beta d_i} \left[\beta \sum_{j = 1}^{\lvert N(i) \rvert} x_j^{\nu} - \nabla f_i(x_i^{\nu}) - \beta \mu_i^{\nu D + \ell}\right]$;
                        
                    $x_i^{\nu+1} = s_i^{\nu} + q_i^{\nu +1}$;
                        
                    \Comment{Every D periods compute dual variable}
                    \eIf{$ \ell < D $}{ 
                        $\mu_i^{\nu D + \ell} = \mu_i^{\nu D}$;
                        
                        $\ell \gets \ell + 1$;
                    }{
                        $\mu_i^{\nu D + \ell} = \mu_i^{\nu D} + \beta \sum_{j = 1}^{\lvert N(i) \rvert} \sum_{\ell = 1}^{\lfloor t/DT \rfloor}(x_i^{\ell} - x_j^{\ell})$;
                            
                        $\ell \gets 1$;
                    }
                }
            }
            $\nu \gets \nu + 1$;
        \end{algorithmic}
    \end{algorithm}

    The new proposed algorithm \ref{alg:reg-prox-gpda} have similar properties as algorithm \ref{alg:prox-gpda}, with only addition the error position variable $q(t)$. The new algorithm have a different optimality condition, which is
    \begin{equation}\label{eq:3.1.12}
        \nabla f(x^{\nu}) + A^T\mu^{\nu D} + \beta A^TAx^{\nu + 1} + \beta B^TB(x^{\nu + 1} - x^{\nu}) - 2\beta \mathrm{D}q^{\nu + 1} = 0
    \end{equation}
    With the optimality condition (\ref{eq:3.1.12}), can be proven the necessary lemmas and theorems in order for the algorithm to have convergence in the optimal point.

\section{Convergence Analysis}

    The new proposed algorithm along with its optimality condition of (\ref{eq:3.1.12}), must reach convergent state in order to be applied in a swarm of agents. The fundamental point of the analysis is the creation of a potential function that decrease at every iteration of the algorithm. The potential function is a synthesis of the augmented Lagrangian and a specific proximal term, along with the magnitude of the violation of the constraint. As a result, it evaluates the effectiveness of both the primary and dual variable updates.

    The first step of the analysis is to provide a bound of the size of the constraint violation using a quantity related to the primal iterates. Let $\underline{\sigma}(A)$ be the minimum singular value of the matrix $A^TA$. Then we have the following result.
    \newpage
    \begin{lemma}\label{lem:3.1}
        Suppose that the assumptions A.\ref{asm:1} and A.\ref{asm:3} are satisfied. Then, the following bound for the dual variable is true for the proposed regulated algorithm.
        
        \begin{equation}\label{eq:3.2.1}
            \begin{align}
                \frac{1}{\beta} \left \lVert \mu^{(\nu + 1)D} - \mu^{\nu D} \right \rVert^2 \leq \frac{3\mathrm{L}^2}{\beta \underline{\sigma}^2(A)} \left \lVert x^{\nu} - x^{\nu - 1}\right \rVert^2 + \frac{12 \beta}{\underline{\sigma}^2(A)} \left \lVert \right\mathrm{D}\left( q^{\nu + 1} - q^{\nu}\right) \rVert^2 \\
                + \frac{3\beta}{\underline{\sigma}^2(A)} \left \lVert B^TB\left( \left(x^{\nu + 1} - x^{\nu}\right) - \left( x^{\nu} - x^{\nu - 1} \right) \right) \right \rVert^2
            \end{align}
        \end{equation}
    \end{lemma}
    \begin{proof}
        In the optimality condition of (\ref{eq:3.1.12}) is applied the (\ref{eq:3.1.10}). So it is
        \begin{equation}\label{eq:3.2.2}
            \nabla f(x^{\nu}) + A^T\mu^{\nu + 1} + \beta B^TB(x^{\nu + 1} - x^{\nu}) - 2\beta \mathrm{D}q^{\nu + 1} = 0
        \end{equation}
        It can be seen that from A.\ref{asm:3}, the vector $0$ lies in the column space of the matrix $A$. Therefore,
        \begin{equation*}
            \mu^{(\nu + 1)D} - \mu^{\nu D} = \beta A x^{\nu + 1} \in \Im(A).
        \end{equation*}
        With that, the difference between the dual variable iteration lies in the column space of the matrix $A$. In addition, from the fact that $\mu^0 = 0$, the dual variable itself lies in the column space of the matrix $A$.
        \begin{equation}\label{eq:3.2.3}
            \mu^{\nu D} = \beta \sum_{\ell = 1}^{\lfloor \nu/D \rfloor} Ax^{\ell}\in \Im(A).
        \end{equation}
        Applying these facts to (\ref{eq:3.2.2}) and $\lambda_{min}(A)$ the smallest nonzero eigenvalue of $A^TA$, it is
        \begin{equation*}
            \lambda_{min}^{1/2}(A) \left \lVert \mu^{(\nu + 1)D} - \mu^{\nu D} \right \rVert \leq \left \lVert A^T \left( \mu^{(\nu + 1)D} - \mu^{\nu D} \right) \right \rVert.
        \end{equation*}
        This inequality, along with equation (\ref{eq:3.2.2}) and $\lambda_{min}^{1/2}(A) = \underline{\sigma}(A)$ implies that
        \begin{equation*}
        \begin{align}
            \left \lVert \mu^{(\nu + 1)D} - \mu^{\nu D} \right \rVert &\leq \frac{1}{\underline{\sigma}(A)}  \left \lVert A^T \left( \mu^{(\nu + 1)D} - \mu^{\nu D} \right) \right \rVert\\
            &\leq \frac{1}{\underline{\sigma}(A)} \left \lVert A^T \mu^{(\nu + 1)D} - A^T\mu^{\nu D} \right \rVert\\
            &\leq \frac{1}{\underline{\sigma}(A)} \left \lVert  -\beta B^TB\left[\left(x^{(\nu + 1)D} - x^{\nu D}\right) - \left(x^{\nu D} - x^{(\nu - 1)D} \right)\right]\right \\
            &\left -\left[\nabla f(x^{\nu D}) - \nabla f(x^{(\nu - 1)D})\right] + 2\beta \mathrm{D}\left(q^{(\nu + 1)D} - q^{\nu D}\right)\right\rVert
            \end{align}
        \end{equation*}
        The next step is to square both sides and divide by $\beta$. Then it is
        \begin{equation*}
            \begin{align}
                \frac{1}{\beta}\left \lVert \mu^{(\nu + 1)D} - \mu^{\nu D} \right \rVert^2 \leq \frac{1}{\beta\underline{\sigma}^2(A)} \left \lVert \beta B^TB\left[ \left(x^{(\nu + 1)D} - x^{\nu D}\right) - \left( x^{\nu D} - x^{(\nu - 1)D} \right)\right] \right\\
                \left+\left[ \nabla f(x^{\nu D}) - \nabla f(x^{(\nu - 1)D})\right] - 2\beta \mathrm{D}\left( q^{(\nu + 1)D} - q^{\nu D}\right)\right\rVert^2
            \end{align}
        \end{equation*}
        By applying A.\ref{asm:1} in the inequality and the trigonal sum inequality, it is
        \begin{equation*}
            \begin{align}
                \frac{1}{\beta}\left \lVert \mu^{(\nu + 1)D} - \mu^{\nu D} \right \rVert^2 \leq \frac{1}{\beta\underline{\sigma}^2(A)}\left[ \beta^2\left \lVert  B^TB\left[ \left(x^{(\nu + 1)D} - x^{\nu D}\right) - \left( x^{\nu D} - x^{(\nu - 1)D} \right)\right]\right\rVert^2\right\\
                +\left\left \lVert\nabla f(x^{\nu D}) - \nabla f(x^{(\nu - 1)D})\right\rVert^2 + 4\beta^2 \left\lVert\mathrm{D}\left( q^{(\nu + 1)D} - q^{\nu D}\right)\right\rVert^2 \right]
            \end{align}
        \end{equation*}
        Now Hölder's inequality, def.\ref{def:2.22}, is applied for $n = 3$. So,
        \begin{equation*}
            \begin{align}
                \frac{1}{\beta}\left \lVert \mu^{(\nu + 1)D} - \mu^{\nu D} \right \rVert^2 \leq \frac{1}{\beta\underline{\sigma}^2(A)}\left[ 3\beta^2\left \lVert  B^TB\left[ \left(x^{(\nu + 1)D} - x^{\nu D}\right) - \left( x^{\nu D} - x^{(\nu - 1)D} \right)\right]\right\rVert^2\right\\
                +\left 3\mathrm{L}^2\left \lVert x^{\nu D} - x^{(\nu - 1)D}\right\rVert^2 + 12\beta^2 \left\lVert\mathrm{D}\left( q^{(\nu + 1)D} - q^{\nu D}\right)\right\rVert^2 \right]    
            \end{align}
        \end{equation*}
        This indicates that the final result of (\ref{eq:3.2.1}) has been proven.
    \end{proof}
    The second step of the analysis is the boundedness of the augmented Lagrangian function. Then the later lemma is proposed.
    \begin{lemma}\label{lem:3.2}
        Suppose that the A.\ref{asm:1} and A.\ref{asm:3} are satisfied. Then the following expression is true for the new proposed algorithm

        \begin{equation}\label{eq:3.2.4}
            \begin{align}
                L_{\beta}(x^{(\nu + 1)D}, \mu^{(\nu + 1)D}) - L_{\beta}(x^{\nu D}, \mu^{\nu D}) \leq -\frac{\gamma - 2\mathrm{L} - \epsilon}{2} \left \lVert x^{(\nu + 1)D} - x^{\nu D}\right \rVert^2 + \frac{2 \beta^2}{\epsilon}\left \lVert \mathrm{D}q^{(\nu+1)D}\right \rVert^2\\ 
                + \frac{3\mathrm{L}^2}{\beta \underline{\sigma}^2(A)}\left \lVert x^{\nu D} - x^{(\nu -1 )D}\right \rVert^2 + \frac{12 \beta}{\underline{\sigma}^2(A)}\left \lVert \mathrm{D}(q^{(\nu+1)D}-q^{\nuD})\right \rVert^2\\
                +\frac{3 \beta}{\underline{\sigma}^2(A)}\left \lVert B^TB\left(\left(x^{(\nu+1)D}-x^{\nu D}\right)-\left(x^{\nu D} - x^{(\nu - 1)D}\right)\right)\right \rVert^2
            \end{align}
        \end{equation}
    \end{lemma}
    \begin{proof}
        The classical Augmented Lagrangian has form
        \begin{equation*}
            L_{\beta}(x^{\nu D}, \mu^{\nu D}) = f(x^{\nu D}) + \langle \mu^{\nu D}, Ax^{\nu D}\rangle + \frac{\beta}{2}\left\lVert Ax^{\nu D}\right\rVert^2
        \end{equation*}
        The potential function of the algorithm is the augmented Lagrangian with the proximal term $\frac{\beta}{2}\left\lVert x - x^{\nu D}\right\rVert^2$. It is known that function $f(x)$ has continuously Lipschitz derivative and \\$A^TA + B^TB = 2\mathrm{D} \succeq 2I$. Then from A.\ref{asm:1} and a result from theorem \ref{thr:2.2} the potential function is strongly convex with modulus $\gamma := 2\beta - \mathrm{L} > 0$. Firstly, the later subtraction will be computed, which is going to be used in the proof of this lemma. In a certain position with different dual variable values, the following expression has been proved.
        \begin{equation*}
            \begin{align}
                L_{\beta}(x^{(\nu + 1)D}, \mu^{(\nu + 1)D})- L_{\beta}(x^{(\nu + 1)D}, \mu^{\nu D}) = f(x^{(\nu + 1)D}) + \langle \mu^{(\nu + 1)D}, Ax^{(\nu + 1)D}\rangle + \frac{\beta}{2}\left\lVert Ax^{(\nu + 1)D}\right\rVert^2 \\
                - f(x^{(\nu + 1)D}) - \langle \mu^{\nu D}, Ax^{(\nu + 1)D}\rangle - \frac{\beta}{2}\left\lVert Ax^{(\nu + 1)D}\right\rVert^2 
            \end{align}
        \end{equation*}

        \begin{equation*}
            \begin{align}
                L_{\beta}(x^{(\nu + 1)D}, \mu^{(\nu + 1)D})- L_{\beta}(x^{(\nu + 1)D}, \mu^{\nu D}) &= \left\langle \mu^{(\nu + 1)D}, Ax^{(\nu + 1)D} \right \rangle -\left\langle \mu^{\nu D}, Ax^{(\nu + 1)D}\right\rangle\\
                &= \left\langle \mu^{(\nu + 1)D} - \mu^{\nu D}, Ax^{(\nu + 1)D}\right\rangle
            \end{align}
        \end{equation*}
        From a result of the previous lemma it is known that the subtraction between two iterations lies into the column space of matrix $A$. So,
        \begin{equation*}
            L_{\beta}(x^{(\nu + 1)D}, \mu^{(\nu + 1)D})- L_{\beta}(x^{(\nu + 1)D}, \mu^{\nu D}) = \left\langle \mu^{(\nu + 1)D} - \mu^{\nu D}, \frac{1}{\beta}\left(\mu^{(\nu + 1)D} - \mu^{\nu D}\right) \right\rangle
        \end{equation*}
        From a property of the inner product, the constant $1/\beta$ will go outside the product. Thus,
        \begin{equation*}
            L_{\beta}(x^{(\nu + 1)D}, \mu^{(\nu + 1)D})- L_{\beta}(x^{(\nu + 1)D}, \mu^{\nu D}) = \frac{1}{\beta}\left\langle \mu^{(\nu + 1)D} - \mu^{\nu D},\mu^{(\nu + 1)D} - \mu^{\nu D}\right\rangle
        \end{equation*}
        Then, the final result will be
        \begin{equation}\label{eq:3.2.5}
            L_{\beta}(x^{(\nu + 1)D}, \mu^{(\nu + 1)D})- L_{\beta}(x^{(\nu + 1)D}, \mu^{\nu D}) = \frac{1}{\beta} \left \lVert \mu^{(\nu + 1)D} - \mu^{\nu D} \right \rVert^2
        \end{equation}
        With the use of the equation (\ref{eq:3.2.5}), the potential function will be:
        \begin{equation*}
            \begin{align}
                L_{\beta}(x^{(\nu + 1)D}, \mu^{(\nu + 1)D}) - L_{\beta}(x^{\nu D}, \mu^{\nu D}) &=  L_{\beta}(x^{(\nu + 1)D}, \mu^{(\nu + 1)D}) - L_{\beta}(x^{(\nu + 1)D}, \mu^{\nu D}) \\&+ L_{\beta}(x^{(\nu + 1)D}, \mu^{\nu D}) - L_{\beta}(x^{\nu D}, \mu^{\nu D}) \\
                &\leq  L_{\beta}(x^{(\nu + 1)D}, \mu^{(\nu + 1)D}) - L_{\beta}(x^{(\nu + 1)D}, \mu^{\nu D}) \\ 
                &+ L_{\beta}(x^{(\nu + 1)D}, \mu^{\nu D})+ \frac{\beta}{2} \left \lVert x^{(\nu + 1)D} - x^{\nu D}\right \rVert_{B^TB}^2\\
                &- L_{\beta}(x^{\nu D}, \mu^{\nu D})  
            \end{align}
        \end{equation*}
        Now is applied equation (\ref{eq:3.2.5}), definition \ref{def:2.25} to the inequality.
        \begin{equation}\label{eq:3.2.6}
            \begin{align}
                L_{\beta}(x^{(\nu + 1)D}, \mu^{(\nu + 1)D}) &\leq L_{\beta}(x^{\nu D}, \mu^{\nu D})\frac{1}{\beta} \left \lVert \mu^{(\nu + 1)D} - \mu^{\nu D} \right \rVert^2 - \frac{\gamma}{2}\left \lVert x^{(\nu + 1)D} - x^{\nu D} \right \rVert^2 \\&+ \left \langle \nabla_x L_{\beta} (x^{(\nu + 1)D}, \mu^{\nu D}) + \beta B^TB(x^{(\nu + 1)D} - x^{\nu D}), x^{(\nu + 1)D} - x^{\nu D}\right \rangle
            \end{align}
        \end{equation}
        From optimality condition of (\ref{eq:3.1.12}) it is
        \begin{equation*}
            \nabla f(x^{\nu D}) + A^T\mu^{\nu D} + \beta A^TAx^{(\nu + 1)D} + \beta B^TB(x^{(\nu + 1)D} - x^{\nu D}) - 2\beta \mathrm{D}q^{(\nu + 1)D} = 0
        \end{equation*}
        \begin{equation*}
            \nabla f(x^{\nu D}) - \nabla f(x^{(\nu + 1)D}) +  \nabla f(x^{(\nu + 1)D}) + A^T\mu^{\nu D} + \beta A^TAx^{(\nu + 1)D} + \beta B^TB(x^{(\nu + 1)D} - x^{\nu D}) - 2\beta \mathrm{D}q^{(\nu + 1)D} = 0
        \end{equation*}
        \begin{equation*}
            \nabla f(x^{\nu D}) - \nabla f(x^{(\nu + 1)D}) + \nabla_x L_{\beta} (x^{(\nu + 1)D}, \mu^{\nu D}) + \beta B^TB(x^{(\nu + 1)D} - x^{\nu D}) - 2\beta \mathrm{D}q^{(\nu + 1)D} = 0
        \end{equation*}
        \begin{equation}\label{eq:3.2.7}
            \nabla_x L_{\beta} (x^{(\nu + 1)D}, \mu^{\nu D}) + \beta B^TB(x^{(\nu + 1)D} - x^{\nu D}) = 2\beta \mathrm{D}q^{(\nu + 1)D} + \nabla f(x^{(\nu + 1)D}) - \nabla f(x^{\nu D})
        \end{equation}
        Now, equation (\ref{eq:3.2.7}) is applied to equation (\ref{eq:3.2.6}). So,
        \begin{equation*}
            \begin{align}
                L_{\beta}(x^{(\nu + 1)D}, \mu^{(\nu + 1)D}) &\leq L_{\beta}(x^{\nu D}, \mu^{\nu D}) + \frac{1}{\beta} \left \lVert \mu^{(\nu + 1)D} - \mu^{\nu D} \right \rVert^2 - \frac{\gamma}{2}\left \lVert x^{(\nu + 1)D} - x^{\nu D} \right \rVert^2 \\
                &+ \left \langle \nabla f(x^{(\nu + 1)D}) - \nabla f(x^{\nu D}), x^{(\nu + 1)D} - x^{\nu D} \right \rangle\\
                &+  \left \langle 2 \beta\mathrm{D}q^{(\nu + 1)D}, x^{(\nu + 1)D} - x^{\nu D} \right \rangle  
            \end{align}
        \end{equation*}
        Assumption A.\ref{asm:1} is applied to the one inner product, so
        \begin{equation}\label{eq:3.2.8}
            \begin{align}
                L_{\beta}(x^{(\nu + 1)D}, \mu^{(\nu + 1)D}) - L_{\beta}(x^{\nu D}, \mu^{\nu D}) &\leq \frac{1}{\beta} \left \lVert \mu^{(\nu + 1)D} - \mu^{\nu D} \right \rVert^2 - \frac{\gamma}{2}\left \lVert x^{(\nu + 1)D} - x^{\nu D} \right \rVert^2 \\
                &+\frac{2\mathrm{L}}{2} \left \lVert x^{(\nu + 1)D} - x^{\nu D} \right \rVert ^2 +  \left \langle 2 \beta \mathrm{D}q^{(\nu + 1)D}, x^{(\nu + 1)D} - x^{\nu D} \right \rangle
            \end{align}
        \end{equation}
        In the second inner product is applied the Young Inequality (definition \ref{def:2.21})
        \begin{equation*}
            \begin{align}
                \left \langle 2 \beta \mathrm{D}q^{(\nu + 1)D}, x^{(\nu + 1)D} - x^{\nu D} \right \rangle \leq \frac{\epsilon}{2}\left \lVert x^{(\nu + 1)D} - x^{\nu D} \right \rVert ^2 + \frac{1}{2\epsilon}\left \lVert 2 \beta \mathrm{D}q^{(\nu + 1)D} \right \rVert ^2\\
                = \frac{\epsilon}{2}\left \lVert x^{(\nu + 1)D} - x^{\nu D} \right \rVert ^2 + \frac{2 \beta^2}{\epsilon} \left \lVert \mathrm{D}q^{(\nu + 1)D} \right \rVert ^2
            \end{align}
        \end{equation*}
        \begin{equation}\label{eq:3.2.9}
            \left \langle 2 \beta \mathrm{D}q^{(\nu + 1)D}, x^{(\nu + 1)D} - x^{\nu D} \right \rangle \leq \frac{\epsilon}{2}\left \lVert x^{(\nu + 1)D} - x^{\nu D} \right \rVert ^2 + \frac{2 \beta^2}{\epsilon} \left \lVert \mathrm{D}q^{(\nu + 1)D} \right \rVert ^2
        \end{equation}
        The equation (\ref{eq:3.2.9}) is applied to equation (\ref{eq:3.2.8}). Therefore,
        \begin{equation}\label{eq:3.2.10}
            \begin{align}
                L_{\beta}(x^{(\nu + 1)D}, \mu^{(\nu + 1)D}) - L_{\beta}(x^{\nu D}, \mu^{\nu D})& \leq \frac{1}{\beta} \left \lVert \mu^{(\nu + 1)D} - \mu^{\nu D} \right \rVert^2 - \frac{\gamma}{2}\left \lVert x^{(\nu + 1)D} - x^{\nu D} \right \rVert^2 \\
                &+\frac{2\mathrm{L}}{2} \left \lVert x^{(\nu + 1)D} - x^{\nu D} \right \rVert ^2 +  \frac{\epsilon}{2}\left \lVert x^{(\nu + 1)D} - x^{\nu D} \right \rVert ^2 \\&+ \frac{2 \beta^2}{\epsilon} \left \lVert \mathrm{D}q^{(\nu + 1)D} \right \rVert ^2
            \end{align}
        \end{equation}
        Lastly, the inequality (\ref{eq:3.2.1}) is applied in the inequality (\ref{eq:3.2.10}).
        \begin{equation}\label{eq:3.2.11}
            \begin{align}
                 L_{\beta}(x^{(\nu + 1)D}, \mu^{(\nu + 1)D}) - L_{\beta}(x^{\nu D}, \mu^{\nu D}) &\leq - \frac{\gamma - 2\mathrm{L} - \epsilon}{2}\left \lVert x^{(\nu + 1)D} - x^{\nu D} \right \rVert^2 + \frac{2 \beta^2}{\epsilon} \left \lVert \mathrm{D}q^{(\nu + 1)D} \right \rVert ^2\\
                 &+\frac{3\mathrm{L}^2}{\beta \underline{\sigma}^2(A)} \left \lVert x^{\nu D} - x^{(\nu - 1)D}\right \rVert^2\\
                 &+ \frac{12 \beta}{\underline{\sigma}^2(A)} \left \lVert \right\mathrm{D}\left( q^{(\nu + 1)D} - q^{\nu D}\right) \rVert^2 \\
                &+ \frac{3\beta}{\underline{\sigma}^2(A)} \left \lVert B^TB\left( \left(x^{(\nu + 1)D} - x^{\nu D}\right) - \left( x^{\nu D} - x^{(\nu - 1)D} \right) \right) \right \rVert^2
            \end{align}
        \end{equation}
        Thus, it is shown the desired inequality.
    \end{proof}

    From Lemma \ref{lem:3.2}, can be observed that the potential function is not decreasing when the iteration increases. This is clearly shown in the terms that have positive sign. In addition, no matter how big value has the constant $\beta$, it does not affect the decrease of the potential function. So, this observation suggests that the classical augmented Lagrangian can not serve as a potential function for the proposed algorithm. Therefore, a new potential Lyapunov function must be defined, in order to prove decrease of the potential function. A new term must be added to the function. The extra term that will be added is
    \begin{equation*}
        \frac{\beta}{2}\left(\left \lVert Ax^{(\nu + 1)D}\right \rVert^2 + \left \lVert x^{(\nu + 1)D} - x^{\nu D} \right \rVert^2_{B^TB} \right)
    \end{equation*}
    This new term will be used to decrease the term $\left \lVert B^TB\left( \left(x^{(\nu + 1)D} - x^{\nu D}\right) - \left( x^{\nu D} - x^{(\nu - 1)D} \right) \right) \right \rVert^2$ after every new iteration. Subsequently, the following lemma will illustrate the descent of the sum of the constraint violation $\left \lVert Ax^{(\nu + 1)D}\right \rVert^2$ and the proximal term $\left \lVert x^{(\nu + 1)D} - x^{\nu D} \right \rVert^2_{B^TB}$ has the desired term.
    
    \begin{lemma}\label{lem:3.3}
        Suppose that assumption A.\ref{asm:1} is satisfied. Then, the following inequality holds true
        \begin{equation}\label{eq:3.2.12}
            \begin{align}
                \frac{\beta}{2}\left(\left \lVert Ax^{(\nu + 1)D}\right \rVert^2 + \left \lVert x^{(\nu + 1)D} - x^{\nu D} \right \rVert^2_{B^TB} \right)
                &\leq \frac{\beta}{2} \left \lVert Ax^{\nu D}\right \rVert^2 + \left \lVert x^{\nu D} - x^{(\nu - 1)D} \right \rVert^2_{B^TB}\\
                &- \frac{\beta}{2}\left \lVert A\left(x^{(\nu + 1)D} - x^{\nu D}\right)\right \rVert^2 
                \\&+ \frac{\mathrm{L} + \epsilon}{2}\left \lVert x^{(\nu + 1)D} - x^{\nu D} \right \rVert^2 + \frac{\mathrm{L}}{2}\left \lVert x^{\nu D} - x^{(\nu - 1)D} \right \rVert^2 \\&+ \frac{2 \beta^2}{\epsilon}\left \lVert \mathrm{D}\left(q^{(\nu + 1)D} - q^{\nu D}\right) \right \rVert^2 \\
                &- \frac{\beta}{2}\left \lVert \left(x^{(\nu + 1)D} -x^{\nu D}\right) - \left( x^{\nu D} - x^{(\nu - 1)D} \right)  \right \rVert^2_{B^TB}
            \end{align}
        \end{equation}
    \end{lemma}
    \begin{proof}
        From the optimality condition of equation (\ref{eq:3.1.12}), it is:
        \begin{equation*}
            \left \langle \nabla f(x^{\nu D}) + A^T\mu^{\nu D} + \beta A^TAx^{(\nu + 1)D} + \beta B^TB(x^{(\nu + 1)D} - x^{\nu D}) - 2\beta \mathrm{D}q^{(\nu + 1)D}, x^{(\nu + 1)D} - x \right \rangle \leq 0, \forall x \in \mathrm{R}^Q
        \end{equation*}
        \begin{equation*}
            \left \langle \nabla f(x^{(\nu - 1)D}) + A^T\mu^{(\nu - 1)D} + \beta A^TAx^{\nu D} + \beta B^TB(x^{\nu D} - x^{(\nu - 1)D}) - 2\beta \mathrm{D}q^{\nu D},  x^{\nu D} - x\right \rangle \leq 0, \forall x \in \mathrm{R}^Q
        \end{equation*}
        Plugging $x = x^{\nu D}$ into the first inequality and $x = x^{(\nu + 1)D}$ into the second, adding the resulting inequalities and utilizing the dual variable update, it is obtained,
        \begin{equation*}
            \begin{align}
                \left \langle  \nabla f(x^{\nu D}) -  \nabla f(x^{(\nu - 1)D}) + A^T(\mu^{(\nu + 1)D} - \mu^{\nu D}) + \beta B^TB\left( \left(x^{\nu D} - x^{(\nu - 1)D} \right) - \left( x^{\nu D} - x^{(\nu - 1)D}\right) \right) \\
                - 2\beta \mathrm{D}(q^{(\nu +1 )D} - q^{\nu D}), x^{(\nu + 1)D} - x^{\nu D}\right \rangle \leq 0
            \end{align}
        \end{equation*}
        Rearranging the above inequality, it is
        \begin{equation}\label{eq:3.2.13}
            \begin{align}
                &\left \langle A^T(\mu^{(\nu + 1)D} - \mu^{\nu D}), x^{(\nu + 1)D} - x^{\nu D} \right \rangle\\& \leq
                - \left \langle \nabla f(x^{\nu D}) -  \nabla f(x^{(\nu - 1)D})  +\beta B^TB\left( \left(x^{\nu D} - x^{(\nu - 1)D} \right) \right) - 2\beta \mathrm{D}(q^{(\nu +1 )D} - q^{\nu D}) ,x^{(\nu + 1)D} - x^{\nu D} \right \rangle
            \end{align}
        \end{equation}
        The next step of the proof is to bound both sides of the inequality. For the lhs of (\ref{eq:3.2.13}) it can be expressed as
        \begin{equation*}
            \begin{align}
                \left \langle A^T(\mu^{(\nu + 1)D} - \mu^{\nu D}), x^{(\nu + 1)D} - x^{\nu D} \right \rangle &= \left \langle \beta A^TAx^{(\nu + 1)D}, x^{(\nu + 1)D} - x^{\nu D} \right \rangle\\
                &= \left \langle \beta Ax^{(\nu + 1)D}, Ax^{(\nu + 1)D} - Ax^{\nu D} \right \rangle\\
                &= \beta \left \lVert Ax^{(\nu + 1)D} \right \rVert^2 - \beta \left \langle Ax^{(\nu + 1)D}, Ax^{\nu D} \right \rangle\\
                &= \frac{\beta}{2} \left( \left \lVert Ax^{(\nu + 1)D} \right \rVert^2 - \left \lVert Ax^{\nu D} \right \rVert^2 + \left \lVert A(x^{(\nu + 1)D} - x^{\nu D}) \right \rVert^2\right)
            \end{align}
        \end{equation*}
        Therefore, the lhs will be:
        \begin{equation} \label{eq:3.2.14}
            \left \langle A^T(\mu^{(\nu + 1)D} - \mu^{\nu D}), x^{(\nu + 1)D} - x^{\nu D} \right \rangle = \frac{\beta}{2} \left( \left \lVert Ax^{(\nu + 1)D} \right \rVert^2 - \left \lVert Ax^{\nu D} \right \rVert^2 + \left \lVert A(x^{(\nu + 1)D} - x^{\nu D}) \right \rVert^2\right)
        \end{equation}
        Secondly, the next step is to bound the rhs of the  inequality (\ref{eq:3.2.13}).
        \begin{equation*}
            \begin{align}
                &- \left \langle \nabla f(x^{\nu D}) -  \nabla f(x^{(\nu - 1)D}) + \beta B^TB\left( \left(x^{\nu D} - x^{(\nu - 1)D} \right) \right) - 2\beta \mathrm{D}(q^{(\nu +1 )D} - q^{\nu D}), x^{(\nu + 1)D} - x^{\nu D} \right \rangle \\
                &= - \left \langle \nabla f(x^{\nu D}) -  \nabla f(x^{(\nu - 1)D}), x^{(\nu + 1)D} - x^{\nu D} \right \rangle - \beta \left \langleB^TB\left( \left(x^{\nu D} - x^{(\nu - 1)D} \right) \right), x^{(\nu + 1)D} - x^{\nu D} \right \rangle\\ &+ 2\beta \left \langle\mathrm{D}(q^{(\nu +1 )D} - q^{\nu D}) ,x^{(\nu + 1)D} - x^{\nu D} \right \rangle\\
                &\leq \mathrm{L} \left \lVert x^{\nu D} - x^{(\nu - 1)D}\right \rVert \left \lVert x^{(\nu + 1)D} - x^{\nu D}\right \rVert - \beta \left \lVert x^{(\nu + 1)D} - x^{\nu D}\right \rVert^2_{B^TB}\\& + \beta \left \langle B^TB( x^{\nu D} - x^{(\nu - 1)D}), x^{(\nu + 1)D} - x^{\nu D}\right \rangle + 2\beta \left \langle\mathrm{D}(q^{(\nu +1 )D} - q^{\nu D}) ,x^{(\nu + 1)D} - x^{\nu D} \right \rangle\\
                &\leq \frac{\mathrm{L}}{2} \left \lVert x^{\nu D} - x^{(\nu - 1)D}\right \rVert^2 + \frac{\mathrm{L}}{2}\left \lVert x^{(\nu + 1)D} - x^{\nu D}\right \rVert - \beta \left \lVert x^{(\nu + 1)D} - x^{\nu D}\right \rVert^2_{B^TB}\\
                 &+ \frac{\beta}{2} \left(\left \lVert x^{\nu D} - x^{(\nu - 1)D}\right \rVert^2_{B^TB} -\left \lVert x^{(\nu + 1)D} - x^{\nu D}\right \rVert^2_{B^TB} - \left \lVert(x^{\nu D} - x^{(\nu - 1)D}) - (x^{(\nu + 1)D} - x^{\nu D})\right \rVert^2_{B^TB}\right)\\&+ \frac{\epsilon}{2}\left \lVert x^{(\nu + 1)D} - x^{\nu D} \right \rVert ^2 + \frac{2 \beta^2}{\epsilon} \left \lVert \mathrm{D}q^{(\nu + 1)D} \right \rVert ^2
            \end{align}
        \end{equation*}
        Thus, the rhs will be:
        \begin{equation}\label{eq:3.2.15}
            \begin{align}
                 &- \left \langle \nabla f(x^{\nu D}) -  \nabla f(x^{(\nu - 1)D}) + \beta B^TB\left( \left(x^{\nu D} - x^{(\nu - 1)D} \right) \right) - 2\beta \mathrm{D}(q^{(\nu +1 )D} - q^{\nu D}), x^{(\nu + 1)D} - x^{\nu D} \right \rangle\\
                 &\leq \frac{\mathrm{L}}{2} \left \lVert x^{\nu D} - x^{(\nu - 1)D}\right \rVert^2 + \frac{\mathrm{L}}{2}\left \lVert x^{(\nu + 1)D} - x^{\nu D}\right \rVert + \frac{\beta}{2}\left \lVert x^{\nu D} - x^{(\nu - 1)D}\right \rVert^2_{B^TB} - \frac{\beta}{2}\left \lVert x^{(\nu + 1)D} - x^{\nu D}\right \rVert^2_{B^TB}\\
                 &- \frac{\beta}{2}\left \lVert(x^{\nu D} - x^{(\nu - 1)D}) - (x^{(\nu + 1)D} - x^{\nu D})\right \rVert^2_{B^TB} + \frac{\epsilon}{2}\left \lVert x^{(\nu + 1)D} - x^{\nu D} \right \rVert ^2 + \frac{2 \beta^2}{\epsilon} \left \lVert \mathrm{D}q^{(\nu + 1)D} \right \rVert ^2
            \end{align}
        \end{equation}
        Combining the inequalities (\ref{eq:3.2.14}) and (\ref{eq:3.2.15}),
        \begin{equation*}
            \begin{align}
                &\frac{\beta}{2} \left( \left \lVert Ax^{(\nu + 1)D} \right \rVert^2 - \left \lVert Ax^{\nu D} \right \rVert^2 + \left \lVert A(x^{(\nu + 1)D} - x^{\nu D}) \right \rVert^2\right)\\
                &\leq \frac{\mathrm{L}}{2} \left \lVert x^{\nu D} - x^{(\nu - 1)D}\right \rVert^2 + \frac{\mathrm{L}}{2}\left \lVert x^{(\nu + 1)D} - x^{\nu D}\right \rVert + \frac{\beta}{2}\left \lVert x^{\nu D} - x^{(\nu - 1)D}\right \rVert^2_{B^TB} - \frac{\beta}{2}\left \lVert x^{(\nu + 1)D} - x^{\nu D}\right \rVert^2_{B^TB}\\
                 &- \frac{\beta}{2}\left \lVert(x^{\nu D} - x^{(\nu - 1)D}) - (x^{(\nu + 1)D} - x^{\nu D})\right \rVert^2_{B^TB} + \frac{\epsilon}{2}\left \lVert x^{(\nu + 1)D} - x^{\nu D} \right \rVert ^2 + \frac{2 \beta^2}{\epsilon} \left \lVert \mathrm{D}q^{(\nu + 1)D} \right \rVert ^2
            \end{align}
        \end{equation*}
        Rearranging the above inequality, the following result has been proven of (\ref{eq:3.2.12}).
        \begin{equation*}
            \begin{align*}
                \frac{\beta}{2}\left(\left \lVert Ax^{(\nu + 1)D}\right \rVert^2 + \left \lVert x^{(\nu + 1)D} - x^{\nu D} \right \rVert^2_{B^TB} \right) &\leq \frac{\beta}{2} \left \lVert Ax^{\nu D}\right \rVert^2 + \frac{\beta}{2}\left \lVert x^{\nu D} - x^{(\nu - 1)D} \right \rVert^2_{B^TB}\\ &- \frac{\beta}{2}\left \lVert A\left(x^{(\nu + 1)D} - x^{\nu D}\right)\right \rVert^2 
                + \frac{\mathrm{L} + \epsilon}{2}\left \lVert x^{(\nu + 1)D} - x^{\nu D} \right \rVert^2\\& + \frac{\mathrm{L}}{2}\left \lVert x^{\nu D} - x^{(\nu - 1)D} \right \rVert^2 \\ &+ \frac{2 \beta^2}{\epsilon}\left \lVert \mathrm{D}\left(q^{(\nu + 1)D} - q^{\nu D}\right) \right \rVert^2\\
                &- \frac{\beta}{2}\left \lVert \left(x^{(\nu + 1)D} -x^{\nu D}\right) - \left( x^{\nu D} - x^{(\nu - 1)D} \right)  \right \rVert^2_{B^TB}
            \end{align*}
        \end{equation*}
    \end{proof}
    
    An interesting fact about the new term of the potential Lyapunov function is that increases in $\frac{\mathrm{L} + \epsilon}{2}\left \lVert x^{(\nu + 1)D} - x^{\nu D} \right \rVert^2$ and $\frac{\mathrm{L}}{2}\left \lVert x^{\nu D} - x^{(\nu - 1)D} \right \rVert^2$, but decreases in\\ $\left \lVert \left(x^{(\nu + 1)D} -x^{\nu D}\right) - \left( x^{\nu D} - x^{(\nu - 1)D} \right)  \right \rVert^2_{B^TB}$. Therefore, the new potential function is difficult to decrease by itself. Fortunately, such extra term can be bounded by the descend of previous terms. Moreover, will be described in the following lemma. Firstly, let a constant $c > 0$. The new potential function for algorithm \ref{alg:reg-prox-gpda} is 
    \begin{equation}\label{eq:3.2.16}
        P_{c, \beta} \left(x^{(\nu + 1)D}, x^{\nu D}, \mu^{(\nu + 1)D}\right) = L_{\beta} \left(x^{(\nu + 1)D}, \mu^{(\nu + 1)D}\right) + \frac{c\beta}{2}\left( \left \lVert Ax^{(\nu + 1)D}\right \rVert^2 + \left \lVert x^{(\nu + 1)D} - x^{\nu D} \right \rVert^2_{B^TB} \right)
    \end{equation}

    \begin{lemma}\label{lem:3.4}
        Suppose that the assumptions made on Lemmas \ref{lem:3.1} and \ref{lem:3.3} are satisfied. Then the following estimation has been proven
        \begin{equation}\label{eq:3.2.17}
            \begin{align}
                P_{c, \beta} \left(x^{(\nu + 1)D}, x^{\nu D}, \mu^{(\nu + 1)D}\right) & \leq  P_{c, \beta} \left(x^{\nu D}, x^{(\nu - 1)D}, \mu^{\nu D}\right) \\
                &- \left(\frac{\gamma - 2\mathrm{L} - \epsilon - c(\mathrm{L} + \epsilon)}{2}\right) \left \lVert x^{(\nu + 1)D} - x^{\nu D}\right\rVert^2\\
                &+ \left( \frac{3\mathrm{L}^2}{\beta \underline{\sigma}^2(A)} + \frac{c\mathrm{L}}{2}\right) \left \lVert x^{\nu D} - x^{(\nu - 1)D}\right\rVert^2\\
                &-\left(\frac{c\beta}{2} - \frac{3\left\lVert B^TB \right \rVert \beta}{\underline{\sigma}^2(A)}\right)\left \lVert \left(x^{(\nu + 1)D} -x^{\nu D}\right) - \left( x^{\nu D} - x^{(\nu - 1)D} \right)  \right \rVert^2_{B^TB}\\
                &+ \frac{2\beta^2}{\epsilon}\left \lVert \mathrm{D} q^{(\nu + 1)D}\right \rVert^2 + \left(\frac{2\beta^2 c}{\epsilon} + \frac{12\beta}{\underline{\sigma}^2(A)}\right) \left \lVert \mathrm{D}( q^{(\nu + 1)D} - q^{\nu D})\right \rVert^2
            \end{align}
        \end{equation}
    \end{lemma}
    \begin{proof}
        In order to prove the decrease of the new potential function, the inequalities (\ref{eq:3.2.11}) and (\ref{eq:3.2.12}) are going to be used. First, suppose a constant $c > 0$ that is going to be multiplied by the second inequality. The next step is to summarize these two inequalities. Therefore,
        \begin{equation*}
            \begin{align}
                &L_{\beta} \left(x^{(\nu + 1)D}, \mu^{(\nu + 1)D}\right)+ \frac{c\beta}{2}\left(\left \lVert Ax^{(\nu + 1)D}\right \rVert^2 + \left \lVert x^{(\nu + 1)D} - x^{\nu D} \right \rVert^2_{B^TB} \right)\\
                &\leq L_{\beta} \left(x^{\nu D}, \mu^{\nu D}\right) +\frac{c\beta}{2} \left \lVert Ax^{\nu D}\right \rVert^2 + \frac{c\beta}{2}\left \lVert x^{\nu D} - x^{(\nu - 1)D} \right \rVert^2_{B^TB} - \frac{c\beta}{2}\left \lVert A\left(x^{(\nu + 1)D} - x^{\nu D}\right)\right \rVert^2 
                \\&+ \frac{c(\mathrm{L} + \epsilon)}{2}\left \lVert x^{(\nu + 1)D} - x^{\nu D} \right \rVert^2 + \frac{\mathrm{cL}}{2}\left \lVert x^{\nu D} - x^{(\nu - 1)D} \right \rVert^2 + \frac{2 \beta^2c}{\epsilon}\left \lVert \mathrm{D}\left(q^{(\nu + 1)D} - q^{\nu D}\right) \right \rVert^2 \\
                &- \frac{\beta}{2}\left \lVert \left(x^{(\nu + 1)D} -x^{\nu D}\right) - \left( x^{\nu D} - x^{(\nu - 1)D} \right)  \right \rVert^2_{B^TB}
                - \frac{\gamma - 2\mathrm{L} - \epsilon}{2}\left \lVert x^{(\nu + 1)D} - x^{\nu D} \right \rVert^2 \\
                &+\frac{3\mathrm{L}^2}{\beta \underline{\sigma}^2(A)} \left \lVert x^{\nu D} - x^{(\nu - 1)D}\right \rVert^2 + \frac{12 \beta}{\underline{\sigma}^2(A)} \left \lVert \right\mathrm{D}\left( q^{(\nu + 1)D} - q^{\nu D}\right) \rVert^2+ \frac{2 \beta^2}{\epsilon} \left \lVert \mathrm{D}q^{(\nu + 1)D} \right \rVert ^2 \\
                &+\frac{3\beta}{\underline{\sigma}^2(A)} \left \lVert B^TB\left( \left(x^{(\nu + 1)D} - x^{\nu D}\right) - \left( x^{\nu D} - x^{(\nu - 1)D} \right) \right) \right \rVert^2
            \end{align}
        \end{equation*}
        applying the definition of the potential function (\ref{eq:3.2.16}), the above inequality will be:
        \begin{equation*}
            \begin{align}
                &P_{c, \beta} \left(x^{(\nu + 1)D}, x^{\nu D}, \mu^{(\nu + 1)D}\right) \leq P_{c, \beta} \left(x^{\nu D}, x^{(\nu - 1)D}, \mu^{\nu D}\right) - \frac{c\beta}{2}\left \lVert A\left(x^{(\nu + 1)D} - x^{\nu D}\right)\right \rVert^2 
                \\&+ \frac{c(\mathrm{L} + \epsilon)}{2}\left \lVert x^{(\nu + 1)D} - x^{\nu D} \right \rVert^2 + \frac{\mathrm{cL}}{2}\left \lVert x^{\nu D} - x^{(\nu - 1)D} \right \rVert^2 + \frac{2 \beta^2c}{\epsilon}\left \lVert \mathrm{D}\left(q^{(\nu + 1)D} - q^{\nu D}\right) \right \rVert^2 \\
                &- \frac{\beta}{2}\left \lVert \left(x^{(\nu + 1)D} -x^{\nu D}\right) - \left( x^{\nu D} - x^{(\nu - 1)D} \right)  \right \rVert^2_{B^TB}
                - \frac{\gamma - 2\mathrm{L} - \epsilon}{2}\left \lVert x^{(\nu + 1)D} - x^{\nu D} \right \rVert^2 \\
                 &+\frac{3\mathrm{L}^2}{\beta \underline{\sigma}^2(A)} \left \lVert x^{\nu D} - x^{(\nu - 1)D}\right \rVert^2 + \frac{12 \beta}{\underline{\sigma}^2(A)} \left \lVert \right\mathrm{D}\left( q^{(\nu + 1)D} - q^{\nu D}\right) \rVert^2 + \frac{2 \beta^2}{\epsilon} \left \lVert \mathrm{D}q^{(\nu + 1)D} \right \rVert ^2 \\
                &+\frac{3\beta}{\underline{\sigma}^2(A)} \left \lVert B^TB\left( \left(x^{(\nu + 1)D} - x^{\nu D}\right) - \left( x^{\nu D} - x^{(\nu - 1)D} \right) \right) \right \rVert^2
            \end{align}
        \end{equation*}
        Therefore, the estimation of the lemma has been proven.
    \end{proof}
    As it can be seen from the previous lemma, there is a positive term in the rhs of the inequality, which involves the term $\left \lVert x^{\nu D} - x^{(\nu - 1)D} \right \rVert^2$. Therefore, the new potential function has a difficulty in decreasing by itself. Clearly, it is observed that such term can be bounded by the descent of the previous iteration. So, taking the summation over all iterations is going to be:
    \begin{equation}\label{eq:3.2.18}
        \begin{align}
            P_{c, \beta} \left(x^{\nu + 1}, x^{\nu}, \mu^{\nu + 1}\right)   
            &\leq P_{c, \beta} \left(x^{1}, x^{0}, \mu^{1}\right) -\left( \frac{\gamma - 2\mathrm{L} - \epsilon - c(\mathrm{L} + \epsilon)}{2} \right)\left \lVert x^{\nu + 1} - x^{\nu} \right \rVert^2\\
             &- \left( \frac{\gamma - 2\mathrm{L} - \epsilon - c(\mathrm{L} + \epsilon)}{2} - \frac{3\mathrm{L}^2}{\beta \underline{\sigma}^2(A)} - \frac{c\mathrm{L}}{2}\right) \sum_{\ell = 1}^{\nu - 1} \left \lVert x^{(\ell + 1)D} - x^{\ell D} \right \rVert^2\\
             &- \left(\frac{c\beta}{2} - \frac{3\beta ||B^TB||}{\underline{\sigma}^2(A)}\right) \sum_{\ell = 1}^{\nu}\left \lVert \left(x^{(\ell + 1)D} -x^{\ell D}\right) - \left( x^{\ell D} - x^{(\ell - 1)D} \right)  \right \rVert^2_{B^TB}
             \\ &+ \frac{2\beta^2}{\epsilon} \sum_{\ell = 1}^{\nu} \left \lVert \mathrm{D}q^{(\ell + 1)D} \right \rVert^2 + \left ( \frac{2\beta^2c}{\epsilon} + \frac{12\beta}{\underline{\sigma}^2(A)}\right) \sum_{\ell = 1}^{\nu} \left \lVert \mathrm{D}(q^{(\ell + 1)D} - q^{\ell D}) \right \rVert^2
        \end{align}
    \end{equation}
    It is clear that the two design parameters must have large values to ensure that the potential function decreases. However, in inequality (\ref{eq:3.2.18}) it is clear that the error variable must be bounded in order to ensure that the potential function is decreasing throughout the iterations. Thus, it must be secured
    \begin{equation}\label{eq:3.2.19}
        \sum_{\ell = 0}^{\infty}\left \lVert q_i^{\ell D}\right\rVert^2 < \infty, \hspace{2.5mm} \sum_{\ell = 0}^{\infty} \left \lVert q_i^{(\ell + 1)D} - q_i^{\ell D} \right \rVert^2 < \infty
    \end{equation}
    The condition (\ref{eq:3.2.19}) will be secured through the design of robust distributed control laws that will ensure that
    \begin{equation}\label{eq:3.2.20}
        \int_{0}^{\infty} \left \lVert q_i(s) \right \rVert ds < \infty, \hspace{5.0mm} \int_{0}^{\infty} \left \lVert q_i(s) \right \rVert^2 ds < \infty
    \end{equation}
    Thus, the variables $q_i(t) \in \mathcal{L}_1 \cap \mathcal{L}_2$.

    \begin{corollary}\label{cor:3.1}
        Suppose that the assumptions A.\ref{asm:1} - A.\ref{asm:3} are satisfied. Then, parameters $c,\beta$ must be chosen such that ensures the decrease of the potential function.
    \end{corollary}
    \begin{proof}
        From the analysis of Lemma \ref{lem:3.4} it is easy to see that as long as $c$ and $\beta$ are chosen large enough, then the potential function will be decreasing at each iteration of the algorithm. So, there be proved lower bounds for these parameters.
        First, it is clear that a sufficient condition for the constant $c$ is derived as:
        \begin{equation*}
            \frac{c\beta}{2} - \frac{3\beta ||B^TB||}{\underline{\sigma}^2(A)} > 0
        \end{equation*}
        \begin{equation*}
            \Rightarrow c > \frac{6||B^TB||}{\underline{\sigma}^2(A)}
        \end{equation*}
        Now, a sufficient condition for the constant $c$ will be 
        \begin{equation}\label{eq:3.2.21}
            c \geq \max \left\{ \frac{\delta}{\mathrm{L}}, \frac{6||B^TB||}{\underline{\sigma}^2(A)} \right\}
        \end{equation}
        As a note, the term $\delta/\mathrm{L}$ in the max operator is going to be used later in the analysis for the convergence of the algorithm. Thus, this is the lower bound of the constant $c$. In addition, the proven bound is independent of the penalty parameter $\beta$.
        Second, for any given constant $c$, the parameter $\beta$ must satisfy the condition
        \begin{equation*}
            \frac{\gamma - 2\mathrm{L} - \epsilon - c(\mathrm{L} + \epsilon)}{2} - \frac{3\mathrm{L}^2}{\beta \underline{\sigma}^2(A)} - \frac{c\mathrm{L}}{2} > 0
        \end{equation*}
        \begin{equation*}
            \Rightarrow \gamma > 2\mathrm{L} + \epsilon +c(\mathrm{L} + \epsilon) + \frac{6\mathrm{L}^2}{\beta\underline{\sigma}^2(A)} + c\mathrm{L}
        \end{equation*}
        \begin{equation*}
            \Rightarrow \beta > \epsilon +c(\mathrm{L} + \epsilon) + \frac{6\mathrm{L}^2}{\beta\underline{\sigma}^2(A)} + c\mathrm{L}
        \end{equation*}
        \begin{equation*}
           \Rightarrow \beta^2 -(2\mathrm{L}c + (1+c)\epsilon)\beta - \frac{6\mathrm{L}^2}{\underline{\sigma}^2(A)} > 0
        \end{equation*}
        So, in order to have a large sufficient value, only the positive root of the above polynomial will be used. Therefore,
        \begin{equation}\label{eq:3.2.22}
            \beta > \frac{1}{2}\left(2\mathrm{L}c + \left(1+c\right)\epsilon + \sqrt{\left(\left(2\mathrm{L}c + \left(1+c\right)\epsilon\right)\right)^2 + \frac{24\mathrm{L}^2}{\underline{\sigma}^2(A)}} \right)
        \end{equation}
    \end{proof}
    It is clear that combining the bounds of $\beta$ and $c$, the result is that $\beta > \delta$. It is concluded that if both bounds (\ref{eq:3.2.21}) and (\ref{eq:3.2.22}) are satisfied, then the potential function $P_{c,\beta}(x^{(\nu + 1)D}, x^{\nu D}, \mu^{(\nu + 1)D})$ decreases at every iteration. With the use of (\ref{eq:3.2.18}) and the particular choice of $c$ and $\beta$, then the constructed potential function will be lower bounded.
    \begin{lemma}\label{lem:3.5}
        Suppose the assumptions A.\ref{asm:1} - A.\ref{asm:3} are satisfied, and $\beta, c$ are chosen according to (\ref{eq:3.2.21}) and (\ref{eq:3.2.22}). Then, if
        \begin{equation*}
            \sum_{\ell = 1}^{T} P_{c,\beta} \left(x^{\ell + 1}, x^{\ell}, \mu^{\ell + 1}\right) > -\infty
        \end{equation*}
        and
        \begin{equation*}
            P_{c,\beta}\left(x^{T + 1}, x^{T}, \mu^{T + 1}\right) \leq P_{c, \beta} \left(x^{T}, x^{T - 1}, \mu^{T}\right) + a^T
        \end{equation*}
        where $a^T \geq 0$ with $\sum_{\ell = 1}^{\infty} a^T < \infty$, then the potential is lower bounded
        \begin{equation}\label{eq:3.2.23}
            \exists \underline{P} \hspace{2.0mm} P_{c, \beta} \left(x^{T}, x^{T - 1}, \mu^{T}\right) \geq \underline{P} > -\infty \hspace{2.0mm} \forall \hspace{2.0mm} T > 0.
        \end{equation}
    \end{lemma}
    \begin{proof}
        To prove lower boundedness, first will be proven the summation of the potential function that is lower bounded. In order to prove it, assumption A.\ref{asm:2} is going to be utilized. First of all the augmented Lagrangian is expressed as follow
        \begin{equation*}
            \begin{align}
                L_{\beta} (x^{(\nu + 1)D},\mu^{(\nu + 1)D}) &= f(x^{(\nu + 1)D}) + \left \langle \mu^{(\nu + 1)D},Ax^{(\nu + 1)D}\right \rangle + \frac{\beta}{2} \left \lVert Ax^{(\nu + 1)D} \right \rVert^2\\
                & = f(x^{(\nu + 1)D}) + \frac{1}{\beta}\left \langle \mu^{(\nu + 1)D},\mu^{(\nu + 1)D} - \mu^{\nu D}\right \rangle + \frac{\beta}{2} \left \lVert Ax^{(\nu + 1)D} \right \rVert^2\\
                & = f(x^{(\nu + 1)D}) + \frac{1}{2\beta}\left ( \left \lVert \mu^{(\nu + 1)D} \right \rVert^2 - \left \lVert \mu^{\nu D}  \right \rVert^2 + \left \lVert \mu^{(\nu + 1)D} - \mu^{\nu D} \right \rVert^2\right) \\&+ \frac{\beta}{2} \left \lVert Ax^{(\nu + 1)D} \right \rVert^2
            \end{align}
        \end{equation*}
        Therefore, summing over $\ell = 1 ... T$, it is obtained
        \begin{equation*}
            \begin{align}
                \sum_{\ell = 1}^{T}  L_{\beta} (x^{(\ell + 1)D},\mu^{(\ell + 1)D})  &= \sum_{\ell = 1}^{T} \left( f(x^{(\ell + 1)D}) + \frac{1}{2\beta}\left \lVert \mu^{(\ell + 1)D} - \mu^{\ell D} \right \rVert^2 + \frac{\beta}{2} \left \lVert Ax^{(\ell + 1)D} \right \rVert^2\right) \\
                &+\frac{1}{2\beta}\left ( \sum_{\ell = 1}^{T}\left \lVert \mu^{(\ell + 1)D} \right \rVert^2 - \sum_{\ell = 1}^{T}\left \lVert \mu^{\ell D}  \right \rVert^2\right)\\ 
                &= \sum_{\ell = 1}^{T} \left( f(x^{(\ell + 1)D}) + \frac{1}{2\beta}\left \lVert \mu^{(\ell + 1)D} - \mu^{\ell D} \right \rVert^2 + \frac{\beta}{2} \left \lVert Ax^{(\ell + 1)D} \right \rVert^2\right) \\
                &+ \frac{1}{2\beta}\left[ \left\lVert \mu^{(T+1)D}\right\rVert^2 + \sum_{\ell = 1}^{T-1}\left\lVert \mu^{(\ell+1)D}\right\rVert^2 - \sum_{\ell = 1}^{T-1}\left\lVert \mu^{\ell D}\right\rVert^2 - \left\lVert \mu^{1}\right\rVert^2\right]\\
                &= \sum_{\ell = 1}^{T} \left( f(x^{(\ell + 1)D}) + \frac{1}{2\beta}\left \lVert \mu^{(\ell + 1)D} - \mu^{\ell D} \right \rVert^2 + \frac{\beta}{2} \left \lVert Ax^{(\ell + 1)D} \right \rVert^2\right) \\
                &+ \frac{1}{2\beta}\left[ \left\lVert \mu^{(T+1)D}\right\rVert^2  - \left\lVert \mu^{1}\right\rVert^2\right]
            \end{align}
        \end{equation*}
       
        Now A.\ref{asm:2} is satisfied and $\beta$ is chosen according to (\ref{eq:3.2.21}) and (\ref{eq:3.2.22}). Then is clear that the above summation is lower bounded since
        \begin{equation*}
            f(x) + \frac{\beta}{2}\left\lVert Ax\right\rVert^2 \geq f(x) + \frac{\delta}{2}\left\lVert Ax\right\rVert^2 \geq 0, \hspace{2.0mm} \forall x \in \mathrm{R}^Q 
        \end{equation*}
        So, the sum of potential function over the $T$ iterations is
        \begin{equation*}
            \begin{align}
                \sum_{\ell = 1}^{T} P_{c,\beta}\left(x^{(\ell + 1)D}, x^{\ell D}, \mu^{(\ell + 1)D}\right) &= \sum_{\ell = 1}^{T} L_{\beta}\left(x^{(\ell + 1)D}, \mu^{(\ell + 1)D}\right) 
                &+\frac{c\beta}{2}\sum_{\ell = 1}^{T}\left(\left \lVert Ax^{(\ell + 1)D}\right \rVert^2 + \left \lVert x^{(\ell + 1)D} - x^{\ell D} \right \rVert^2_{B^TB}\right)\\
                &= \sum_{\ell = 1}^{T} \left( f(x^{(\ell + 1)D}) + \frac{1}{2\beta}\left \lVert \mu^{(\ell + 1)D} - \mu^{\ell D} \right \rVert^2 + \frac{\beta}{2} \left \lVert Ax^{(\ell + 1)D} \right \rVert^2\right) \\
                &+ \frac{1}{2\beta}\left[ \left\lVert \mu^{(T+1)D}\right\rVert^2  - \left\lVert \mu^{1}\right\rVert^2\right]+ \frac{c\beta}{2}\sum_{\ell = 1}^{T}\left(\left \lVert Ax^{(\ell + 1)D}\right \rVert^2 + \left \lVert x^{(\ell + 1)D} - x^{\ell D} \right \rVert^2_{B^TB}\right)\\
                &\geq \sum_{\ell = 1}^{T} \left( f(x^{(\ell + 1)D}) + \frac{\delta}{2} \left \lVert Ax^{(\ell + 1)D} \right \rVert^2 + \frac{1}{2\beta}\left \lVert \mu^{(\ell + 1)D} - \mu^{\ell D} \right \rVert^2\right) \\
                &+ \frac{1}{2\beta}\left[ \left\lVert \mu^{(T+1)D}\right\rVert^2  - \left\lVert \mu^{1}\right\rVert^2\right] \\
                &+ \frac{c\beta}{2}\sum_{\ell = 1}^{T}\left(\left \lVert Ax^{(\ell + 1)D}\right \rVert^2 + \left \lVert x^{(\ell + 1)D} - x^{\ell D} \right \rVert^2_{B^TB}\right) \geq 0
            \end{align}
        \end{equation*}
        \begin{equation*}
            \Rightarrow \sum_{\ell = 1}^{T} P_{c,\beta}\left(x^{(\ell + 1)D}, x^{\ell D}, \mu^{(\ell + 1)D}\right) > -\infty, \hspace{2.0mm} \forall T > 0.
        \end{equation*}
        The summation of the potential function has been proven that is lower bounded and positive. The next step is to ensure that the potential function is upper and lower bounded. For more iterations it is known that
        \begin{equation*}
            \begin{align}
                -\infty &< \sum_{\ell = 1}^{RT} P_{c,\beta}\left(x^{(\ell + 1)D}, x^{\ell D}, \mu^{(\ell + 1)D}\right)\\
                &= P_{c,\beta}\left(x^{ 1}, x^{0}, \mu^{1}\right) + P_{c,\beta}\left(x^{2}, x^{1}, \mu^{2}\right) + ...
                + P_{c,\beta}\left(x^{(RT + 1)D}, x^{RT D}, \mu^{(RT + 1)D}\right)\\
                &\leq \sum_{\ell = 1}^{T-1} P_{c,\beta}\left(x^{(\ell + 1)D}, x^{\ell D}, \mu^{(\ell + 1)D}\right) + P_{c,\beta}\left(x^{TD}, x^{(T-1) D}, \mu^{T D}\right) + P_{c,\beta}\left(x^{(T + 1)D}, x^{T D}, \mu^{(T +1)D}\right) \\&+...+ P_{c,\beta}\left(x^{(RT + 1)D}, x^{RT D}, \mu^{(RT + 1)D}\right)\\
                &\leq \sum_{\ell = 1}^{T-1} P_{c,\beta}\left(x^{(\ell + 1)D}, x^{\ell D}, \mu^{(\ell + 1)D}\right) +\underbrace{P_{c,\beta}\left(x^{TD}, x^{(T-1) D}, \mu^{T D}\right)+...+P_{c,\beta}\left(x^{TD}, x^{(T-1) D}, \mu^{T D}\right)}_\text{R terms}\\
                & + a^T + 2a^T + ... + RTa^T\\
                &= \sum_{\ell = 1}^{T-1} P_{c,\beta}\left(x^{(\ell + 1)D}, x^{\ell D}, \mu^{(\ell + 1)D}\right) + RP_{c,\beta}\left(x^{TD}, x^{(T-1) D}, \mu^{T D}\right) + a^T\sum_{\ell = 1}^{RT}\ell \\
                &=\sum_{\ell = 1}^{T - 1} P_{c,\beta}\left(x^{(\ell + 1)D}, x^{\ell D}, \mu^{(\ell + 1)D}\right) + RP_{c,\beta}\left(x^{TD}, x^{(T-1) D}, \mu^{T D}\right) + \frac{RT(RT + 1)}{2}a^T
            \end{align}
        \end{equation*}
        Hence, both the sum and $a^T$ are lower bounded and positive, then and the potential function is lower bounded. Thus, there exist a low bound for the potential function. In addition, with the proper selection of the parameters $c,\beta$ the potential function is nonincreasing. Therefore, the desired bound has been proven.
    \end{proof}
    
    All the above results are going to be used for the basic theorem of this thesis. The next and final step is to evaluate the convergence of the algorithm and the rate of convergence. Let $Q(x^{(\nu + 1)D}, \mu^{(\nu + 1)D})$ be the optimality gap of problem (\ref{eq:2.5.6}), which is given by
    \begin{equation}\label{eq:3.2.24}
        Q(x^{(\nu + 1)D}, \mu^{(\nu + 1)D}) := \left \lVert \nabla_x L_{\beta} (x^{(\nu + 1)D},\mu^{(\nu + 1)D}) \right \rVert^2 + \beta\left \lVert Ax^{(\nu + 1)D} \right \rVert^2
    \end{equation}
    With equation (\ref{eq:3.2.24}) is clear that when it tends to zero, with the limit be the point $(x_{\nu}^{\bullet}, \mu_{\nu}^{\bullet})$, it implies that it is KKT of problem (\ref{eq:2.5.6}). That means
    \begin{equation}\label{eq:3.2.25}
        \begin{align}
            \nabla f(x_{\nu}^{\bullet}) + A^T \mu_{\nu}^{\bullet} &= 0\\
            Ax_{\nu}^{\bullet} &= 0\\
            I^TA &= 0
        \end{align}
    \end{equation}
    By the use of these equations it can be shown that
    \begin{equation*}
        I^T \nabla f(x_{\nu}^{\bullet}) = 0 \Rightarrow
        \sum_{i = 1}^{N} \nabla f_i(x_{\nu}^{\bullet}) = 0
    \end{equation*}
    In order to reach the optimal condition, the consensus must be achieved, i.e. $Ax^{(\nu + 1)D} \to 0$. Therefore, a necessary condition is that the sequence of the dual variable must converge to 0, which is achieved when the sequence of the primal variable converges. The first condition of (\ref{eq:3.2.25}) can be proved from the optimality gap when it tends to zero. Therefore,
    \begin{equation*}
        \lim_{\nu \to \infty} \nabla f(x^{\nu D}) + A^T\mu^{\nu D} + \beta A^TAx^{(\nu + 1)D} = \nabla f(x_{\nu}^{\bullet}) + A^T \mu_{\nu}^{\bullet} = 0
    \end{equation*}
    The next and final step of the analysis is to ensure the convergence of the algorithm through the optimality gap $Q(\cdot)$. The optimality gap tends to zero in a sublinear manner. This is presented in the following theorem about the proposed algorithm.
    \begin{theorem}\label{thr:3.1}
        Suppose that the main assumptions A.\ref{asm:1} - A.\ref{asm:4} are true and satisfied. In addition, suppose that the parameters $c$ and $\beta$ are chosen such as are presented in (\ref{eq:3.2.21}) and (\ref{eq:3.2.22}), respectively. Then, the following claims have been proven for the sequence which is generated by the proposed algorithm.
        \begin{itemize}
            \item The constraint is satisfied in the limit, i.e.
                \begin{equation*}
                    \lim_{\nu \to \infty} \mu^{(\nu + 1)D} - \mu^{\nu D} \to 0
                \end{equation*}
                \begin{equation*}
                    \lim_{\nu \to \infty} Ax^{(\nu + 1)D} \to 0
                \end{equation*}
                \begin{equation*}
                    \lim_{\nu \to \infty} x^{(\nu + 1)D} - x^{\nu D} \to 0
                \end{equation*}
                \begin{equation*}
                    \lim_{t \to \infty} x(t) \to x^*
                \end{equation*}
            \item  Further, if $\left \lVert f(x) \rVert \right$ is bounded for all $x \in \mathrm{R}^Q$, then the sequence $\{\mu^{\ell D}\}_{\ell = 1}^{\nu + 1}$ is also bounded; If $f(x) + \frac{\beta}{2}\left \lVert Ax \rVert \right^2$ is coercive, then the sequence $\{x^{\ell D}\}_{\ell = 1}^{\nu + 1}$ is bounded.
            \item  Every limit point of the iterates $\{x^{\ell D},\mu^{\ell D}\}_{\ell = 1}^{\nu + 1}$ generated by Algorithm \ref{alg:reg-prox-gpda} converges to a stationary point of problem (\ref{eq:2.5.6}). In addition, $Q(x^{(\nu +1)D}\,\mu^{\nu D}) \to 0$.
            \item For any given $\phi > 0$, let us define $Y$ to be the first that the optimality gap reaches below $\phi$, i.e.,
                \begin{equation*}
                    Y := arg \min_{\nu} Q(x^{(\nu +1)D},\mu^{\nu D}) \leq \phi
                \end{equation*}
            Then there exists a constant $v > 0$ such that the following inequality is true
                \begin{equation*}
                    \phi \leq \frac{v}{Y - 1}
                \end{equation*}
            That is, the optimality gap $Q(x^{(\nu +1)D},\mu^{\nu D})$ converges sublinearly.
        \end{itemize}
    \end{theorem}
    \newpage
    \begin{proof}
        \hspace{5.0mm}
        \begin{enumerate}
            \item The first step is to ensure that the constraint is satisfied in the limit, while the primal variable converges to the optimal point. In order to achieve this, lemmas \ref{lem:3.4} and \ref{lem:3.5} are going to be used. In addition, the NOCGPROX variables must be square integrable. From the results of both lemmas, it is known that in every iteration the potential function decreases and is lower bounded. Thus, it is concluded that $\left \lVert x^{(\nu + 1)D} - x^{\nu D} \right \rVert ^2 \to 0$ and $\left \lVert \left(x^{(\nu + 1)D} -x^{\nu D}\right) - \left( x^{\nu D} - x^{(\nu - 1)D} \right)  \right \rVert^2_{B^TB} \to 0$. Therefore, $x^{(\nu + 1)D} \to x^{\nu D}$ and $x^{(\nu + 1)D} - x^{\nu D} \to x^{\nu D} - x^{(\nu - 1)D}$. In addition, with the use of lemma \ref{lem:3.1} it is concluded that $\left \lVert \mu^{(\nu + 1)D} - \mu^{\nu D} \right \rVert^2 \to 0$. So the dual variable also converges, $\mu^{(\nu + 1)D} \to \mu^{\nu D}$. Hence it is known that the subtraction between the consecutive iterations belong in the $\Im(A)$, then the constraint is satisfied in the limit, $Ax^{(\nu + 1)D} \to 0$, and consensus between the agents has been achieved. So, the agents will converge in the neighborhood of the optimal solution.
            
            \item The optimality condition of algorithm \ref{alg:reg-prox-gpda} is
            \begin{equation*}
                \nabla f(x^{\nu D}) + A^T\mu^{\nu D} + \beta A^TAx^{(\nu + 1)D} + \beta B^TB(x^{(\nu + 1)D} -x^{\nu D}) -2\beta\mathrm{D}q^{(\nu + 1)D} = 0
            \end{equation*}
            As a first step, the sequence $\{\mu^{\ell D}\}_{\ell = 1}^{\nu + 1}$ must be bounded, if $\nabla f(x^{\nu D})$ is bounded. This statement is true, hence it is known from step 1 that $\left \lVert x^{(\nu + 1)D} - x^{\nu D} \right \rVert^2 \to 0$ and $Ax^{(\nu + 1)D} \to 0$. This implies that both the subtraction and the constraint are bounded. Then the boundedness of the dual variable follows from the assumption that $\nabla f(x)$ is bounded for any $x \in \mathbf{R}^Q$ and that the variable lies inside the column space of the incidence matrix of the system. Therefore, the sequence $\{\mu^{\ell D}\}_{\ell = 1}^{\nu + 1}$ is bounded.
            
            The nest step is to ensure boundedness for the primal variable if $f(x) + \frac{\beta}{2} \left \lVert Ax \right \rVert^2$ is coercive. The potential function is
            \begin{equation*}
                \begin{align}
                    P_{c, \beta} (x^{(\nu + 1)D}, x^{\nu D}, \mu^{(\nu + 1)D}) = &f(x^{(\nu + 1)D}) + \left \langle \mu^{(\nu + 1)D}, Ax^{(\nu + 1)D}\right\rangle + \frac{\beta}{2} \left \lVert Ax^{(\nu + 1)D} \right \rVert^2\\
                    &+ \frac{c\beta}{2}\left(\left \lVert Ax^{(\nu + 1)D} \right \rVert^2 + \left \lVert x^{(\nu + 1)D} - x^{\nu D}\right \rVert^2_{B^TB} \right)
                \end{align}
            \end{equation*}
            \begin{equation*}
                \begin{align}
                    \Rightarrow P_{c, \beta} (x^{(\nu + 1)D}, x^{\nu D}, \mu^{(\nu + 1)D}) = &f(x^{(\nu + 1)D}) + \frac{\beta}{2} \left \lVert Ax^{(\nu + 1)D} \right \rVert^2\\
                     &+ \frac{1}{2\beta}\left( \left\lVert \mu^{(\nu + 1)D}\right\rVert^2 - \left\lVert \mu^{\nu D}\right\rVert^2 + \left\lVert \mu^{(\nu + 1)D} - \mu^{\nu D}\right\rVert^2\right)\\
                    &+ \frac{c\beta}{2}\left(\left \lVert Ax^{(\nu + 1)D} \right \rVert^2 + \left \lVert x^{(\nu + 1)D} - x^{\nu D}\right \rVert^2_{B^TB} \right)
                \end{align}
            \end{equation*}
            It is known that from lemma \ref{lem:3.5} that the potential function is decreasing, so it has an upper bound. Suppose that the sequence $\{x^{\ell D}\}_{\ell = 1}^{\nu + 1}$ is not bounded and a domain $\mathcal{K}$ that is an infinite subset of the iteration index in which $\lim_{\nu \in \mathcal{K}}x^{\nu D} = \infty$. Now, this limit is going to be passed to the potential function. So, in this infinite subset the potential function has limit
            \begin{equation*}
                \lim_{\nu \in \mathcal{K}} P_{c\beta}(x^{(\nu + 1)D}, x^{\nu D}, \mu^{(\nu + 1)D}) = \lim_{\nu \in \mathcal{K}} f(x^{(\nu + 1)D}) + \frac{c\beta + \beta}{2}\left \lVert Ax^{(\nu + 1)D} \right \rVert^2 = \infty
            \end{equation*}
            Some terms are zero, hence $\mu^{(\mu + 1)D} \to \mu^{\nu D}$ and $x^{(\nu + 1)D} \to x^{\nu D}$. In addition, the last result of the limit comes from the coerciveness assumption. However, this contradicts the fact that the potential function is upper bounded. So, the sequence $\{x^{\ell D}\}_{\ell = 1}^{\nu + 1}$ is bounded.

            \item Let $\mathcal{K}$ be any converging infinite iteration index, such that $\{(\mu^{\nu D}, x^{\nu D})\}_{\nu \in \mathcal{K}}$ converges to the limit point $(x_{\nu}^{\bullet}, \mu_{\nu}^{\bullet})$. Then, passing the limit in $\mathcal{K}$, and using the fact that the subtraction between two consecutive iterations converge to 0,
            \begin{equation*}
                \nabla f(x_{\nu}^{\bullet}) + A^T\mu_{\nu}^{\bullet} + \beta A^TAx_{\nu}^{\bullet} = 0
            \end{equation*}
            Also, $Ax_{\nu}^{\bullet} = 0$, so it is concluded that $(x_{\nu}^{\bullet}, \mu_{\nu}^{\bullet})$ is a stationary point of the original problem (\ref{eq:2.5.6}). In addition, even if the sequence $\{(\mu^{\nu D}, x^{\nu D})\}_{\ell = 1}^{\nu + 1}$ does not have a limit point, then from bullet 1, is still bounded and converges, i.e. $\left \lVert \mu^{(\nu + 1)D} - \mu^{\nu D} \right \rVert \to 0$, $\left \lVert x^{(\nu + 1)D} - x^{\nu D} \right \rVert \to 0$ and $\left \lVert q^{(\nu + 1)D} \right \rVert \to 0$. Therefore,
            \begin{equation*}
                \begin{align}
                    \lim_{\nu \to \infty} \nabla_x L_{\beta}(x^{(\nu + 1)D}, \mu^{\nu D}) & = \lim_{\nu \to \infty} \nabla f(x^{(\nu + 1)D}) + A^T\mu^{(\nu + 1)D} + A^T(\mu^{\nu D} - \mu^{(\nu + 1)D})\\
                    & = \lim_{\nu \to \infty} -\beta B^TB(x^{(\nu + 1)D} - x^{\nu D}) + A^T(\mu^{\nu D} - \mu^{(\nu + 1)D}) + 2\beta\mathrm{D}q^{(\nu + 1)D}\\ 
                    & = 0
                \end{align}
            \end{equation*}
            where the change of the terms are a result from the optimality condition in addition to the results of step 1 for the convergence. Thus, $Q(x^{(\nu + 1)D}, \mu^{\nu D}) \to 0$.

            \item The first step is to bound the size of the augmented Lagrangian. From the optimality condition of Algorithm \ref{alg:reg-prox-gpda}, it is
            \begin{equation*}
                \begin{align}
                    \left \lVert \nabla_x L_{\beta} (x^{\nu D}, \mu^{(\nu - 1)D}) \right \rVert^2 & = \left \lVert \nabla_x L_{\beta}(x^{(\nu + 1)D}, \mu^{\nu D}) + \beta B^TB(x^{(\nu + 1)D} - x^{\nu D}) \right\\
                    & \left - 2\beta\mathrm{D}q^{(\nu + 1)D} - \nabla_x L_{\beta}(x^{\nu D}, \mu^{(\nu - 1)D}) \right \rVert^2\\
                    & = \left \lVert\nabla f(x^{(\nu + 1)D}) - \nabla f(x^{\nu D})  + A^T(\mu^{(\nu + 1)D} - \mu^{\nu D}) \right \\
                    &\left+ \beta B^TB(x^{(\nu + 1)D} - x^{\nu D})- 2\beta\mathrm{D}q^{(\nu + 1)D} \right \rVert^2\\
                    & \leq \left(\left \lVert \nabla f(x^{(\nu + 1)D}) - \nabla f(x^{\nu D}) \right \rVert + \left \lVert A^T(\mu^{(\nu + 1)D} - \mu^{\nu D}) \right \rVert \right\\
                    &\left + \beta\left \lVert B^TB(x^{(\nu + 1)D} - x^{\nu D}) \right \rVert + 2\beta \left \lVert \mathrm{D}q^{(\nu + 1)D}\right \rVert \right)^2\\
                    & \leq \left( \mathrm{L}\left \lVert x^{(\nu + 1)D} - x^{\nu D} \right \rVert + \left \lVert A^T(\mu^{(\nu + 1)D} - \mu^{\nu D}) \right \rVert \right\\
                    &\left + \beta\left \lVert B^TB(x^{(\nu + 1)D} - x^{\nu D}) \right \rVert + 2\beta \left \lVert \mathrm{D}q^{(\nu + 1)D}\right \rVert \right)^2
                \end{align}
            \end{equation*}
            Now from Holder's inequality, definition \ref{def:2.22}, for $n = 3$, is.
            \begin{equation*}
                \begin{align}
                    \left \lVert \nabla_x L_{\beta} (x^{\nu D}, \mu^{(\nu - 1)D}) \right \rVert^2 & \leq 3\mathrm{L}^2\left \lVert x^{(\nu + 1)D} - x^{\nu D} \right \rVert^2 + 3\left \lVert A^T(\mu^{(\nu + 1)D} - \mu^{\nu D})\right \rVert^2\\
                    &+ 3\beta^2\left \lVert B^TB(x^{(\nu + 1)D} - x^{\nu D})\right \rVert^2 + 12\beta^2 \left \lVert \mathrm{D}q^{(\nu + 1)D}\right \rVert^2\\
                    & \leq 3\mathrm{L}^2\left \lVert x^{(\nu + 1)D} - x^{\nu D} \right \rVert^2 + 3\left \lVert \mu^{(\nu + 1)D} - \mu^{\nu D}\right \rVert^2\left\lVert A^TA\right\rVert\\
                    &+ 3\beta^2\left \lVert B^TB(x^{(\nu + 1)D} - x^{\nu D})\right \rVert^2 + 12\beta^2 \left \lVert \mathrm{D}q^{(\nu + 1)D}\right \rVert^2
                \end{align}
            \end{equation*}
            By utilizing lemma \ref{lem:3.1}, it will be applied into the optimality gap, in order to bound it as well. So,
            \begin{equation*}
                \begin{align}
                    Q(x^{\nu D}, \mu^{(\nu - 1)D}) &= \left \lVert\nabla_x L_{\beta} (x^{\nu D}, \mu^{(\nu - 1)D})\right \rVert^2 + \beta \left \lVert Ax^{\nu D} \right \rVert^2\\
                    & \leq 3\mathrm{L}^2\left \lVert x^{(\nu + 1)D} - x^{\nu D} \right \rVert^2 + 3\left \lVert \mu^{(\nu + 1)D} - \mu^{\nu D}\right \rVert^2\left\lVert A^TA\right\rVert\\
                    & + 3\beta^2\left \lVert B^TB(x^{(\nu + 1)D} - x^{\nu D})\right \rVert^2 + 12\beta^2 \left \lVert \mathrm{D}q^{(\nu + 1)D}\right \rVert^2 + \beta \left \lVert Ax^{\nu D} \right \rVert^2\\
                    & \leq 3\mathrm{L}^2\left \lVert x^{(\nu + 1)D} - x^{\nu D} \right \rVert^2 + 3||A^TA||\left(\frac{3\beta\mathrm{L}^2}{\underline{\sigma}^2(A)} \left \lVert x^{\nu D} - x^{(\nu - 1)D}\right \rVert^2\right\\
                    &+ \frac{12 \beta^2}{\underline{\sigma}^2(A)} \left \lVert \right\mathrm{D}\left( q^{(\nu + 1)D} - q^{\nu D}\right) \rVert^2 \\
                    &\left+ \frac{3\beta^2}{\underline{\sigma}^2(A)} \left \lVert B^TB\left( \left(x^{(\nu + 1)D} - x^{\nu D}\right) - \left( x^{\nu D} - x^{(\nu - 1)D} \right) \right) \right \rVert^2\right)\\
                    & + 3\beta^2\left \lVert B^TB(x^{(\nu + 1)D} - x^{\nu D})\right \rVert^2 + 12\beta^2 \left \lVert \mathrm{D}q^{(\nu + 1)D}\right \rVert^2 + \beta \left \lVert Ax^{\nu D} \right \rVert^2\\
                    & = 3\mathrm{L}^2\left \lVert x^{(\nu + 1)D} - x^{\nu D} \right \rVert^2 + \frac{9\beta\mathrm{L}^2 ||A^TA||}{\underline{\sigma}^2(A)} \left \lVert x^{\nu D} - x^{(\nu - 1)D}\right \rVert^2\\
                    & + \frac{36 \beta^2 ||A^TA||}{\underline{\sigma}^2(A)} \left \lVert \right\mathrm{D}\left( q^{(\nu + 1)D} - q^{\nu D}\right) \rVert^2 \\
                    & + \frac{9\beta^2 ||A^TA||}{\underline{\sigma}^2(A)} \left \lVert B^TB\left( \left(x^{(\nu + 1)D} - x^{\nu D}\right) - \left( x^{\nu D} - x^{(\nu - 1)D} \right) \right) \right \rVert^2\\
                    & + 3\beta^2\left \lVert B^TB(x^{(\nu + 1)D} - x^{\nu D})\right \rVert^2 + 12\beta^2 \left \lVert \mathrm{D}q^{(\nu + 1)D}\right \rVert^2 + \beta \left \lVert Ax^{\nu D} \right \rVert^2\\
                \end{align}
            \end{equation*}
            As it is shown, the optimality gap is upper bounded. So there exists $\chi_1,\chi_2 > 0$ such that the optimality gap is bounded by
            \begin{equation*}
                Q(x^{\nu D}, \mu^{(\nu - 1)D}) \leq \chi_1\left \lVert x^{(\nu + 1)D} - x^{\nu D}\right \rVert^2 + \chi_2\left \lVert B^TB\left( \left(x^{(\nu + 1)D} - x^{\nu D}\right) - \left( x^{\nu D} - x^{(\nu - 1)D} \right) \right) \right \rVert^2 
            \end{equation*}
            From the decent estimate of the potential function, there also exists two constants $\psi_1, \psi_2 > 0$ such that the descent is lower bounded and negative
            \begin{equation*}
                \begin{align}
                    &P_{c,\beta}(x^{(\nu + 1)D}, x^{\nu D}, \mu^{(\nu + 1)D}) - P_{c,\beta} (x^{\nu D}, x^{(\nu - 1)D}, \mu^{\nu D}) \\
                    &\leq -\psi_1 \left \lVert x^{(\nu + 1)D} - x^{\nu D}\right \rVert^2 - \psi_2 \left \lVert B^TB\left( \left(x^{(\nu + 1)D} - x^{\nu D}\right) - \left( x^{\nu D} - x^{(\nu - 1)D} \right) \right) \right \rVert^2 
                \end{align}
            \end{equation*}
            So, matching the above conditions, the optimality gap will be bounded as
            \begin{equation*}
                Q(x^{\nu D}, \mu^{(\nu - 1)D}) \leq \frac{\min\{\psi_1,\psi_2\}}{\max\{\chi_1,\chi_2\}} \left(P_{c,\beta} (x^{\nu D}, x^{(\nu - 1)D}, \mu^{\nu D}) - P_{c,\beta}(x^{(\nu + 1)D}, x^{\nu D}, \mu^{(\nu + 1)D})\right)
            \end{equation*}
            Now, summing over $\ell$ iterations and let $Y$ denote the first time that the optimality gap reaches below $\phi$, it is obtained
            \begin{equation*}
                \begin{align}
                    \phi &\leq \frac{1}{Y - 1} \sum_{\ell = 1}^{Y} Q(x^{\ell D}, \mu^{(\ell - 1)D})\\
                    &\leq \frac{1}{Y - 1}\frac{\min\{\psi_1,\psi_2\}}{\max\{\chi_1,\chi_2\}}\sum_{\ell = 1}^{Y}\left(P_{c,\beta} (x^{\ell D}, x^{(\ell - 1)D}, \mu^{\ell D}) - P_{c,\beta}(x^{(\ell + 1)D}, x^{\ell D}, \mu^{(\ell + 1)D})\right)\\
                    & \leq  \frac{1}{Y - 1}\frac{\min\{\psi_1,\psi_2\}}{\max\{\chi_1,\chi_2\}} \left(P_{c,\beta} (x^1, x^0, \mu^1) - P_{c,\beta}(x^{(Y + 1)D}, x^{Y D}, \mu^{(Y + 1)D})\right)\\
                    & \leq \frac{1}{Y - 1}\frac{\min\{\psi_1,\psi_2\}}{\max\{\chi_1,\chi_2\}} \left(P_{c,\beta} (x^1, x^0, \mu^1) - (P_{c,\beta}(x^{Y D}, x^{(Y - 1)D}, \mu^{Y D}) + a^T)\right)\\
                    & \leq...\leq \frac{1}{Y - 1}\frac{\min\{\psi_1,\psi_2\}}{\max\{\chi_1,\chi_2\}} \left(P_{c,\beta} (x^1, x^0, \mu^1) - P_{c,\beta}(x^2, x^1, \mu^2) - Ya^T\right)\\
                    & \leq \frac{1}{Y - 1}\frac{\min\{\psi_1,\psi_2\}}{\max\{\chi_1,\chi_2\}} \left(P_{c,\beta} (x^1, x^0, \mu^1) - \underline{P}\right)\\
                    &:=\frac{v}{Y - 1}
                \end{align}
            \end{equation*}
            As a conclusion, it has been proven that the optimality gap function $Q(\cdot)$ converges in a sublinear manner.
        \end{enumerate}
    \end{proof}

    With the results of theorem \ref{thr:3.1} it is proven that the discrete form of the algorithm can be applied in a real network of agents. However, depended on the properties of the cost functions the sequence $\{(\mu^{\nu D}, x^{\nu D})\}_{\ell = 1}^{\nu + 1}$ may or may not be bounded. On the contrary, the optimality gap converges to zero in a sublinear manner, thus the agents will always converge to a stationary optimal point. In order to have convergence, the proposed variables must also converge to zero, and be square integrable. So, there must be design specific robust control laws that ensures that the variables tend to zero. 